% !TEX program = xelatex
\documentclass[11pt,a4paper,twoside]{report}

\usepackage{xcolor}
\usepackage[a4paper,margin=28mm]{geometry}
\usepackage{graphicx}
\usepackage{microtype}
\usepackage{setspace}
\usepackage{parskip}
\usepackage{booktabs}
\usepackage{tabularx}
\usepackage{longtable}
\usepackage{array}
\usepackage{multirow}
\usepackage{enumitem}
\usepackage{amssymb}
\usepackage{listings}
\usepackage[hyphens]{url}
\usepackage{tcolorbox}
\usepackage{tikz}
\usepackage[table]{colortbl}
\usetikzlibrary{arrows.meta,positioning,shapes.geometric,fit,calc,chains,
decorations.pathmorphing,backgrounds}

\newtcolorbox{adrbox}[2][]{colback=blue!5,colframe=blue!75!black,fonttitle=\bfseries,title=ADR: #2,#1}

% --- Semantic Hyperlinking --------------------------------------------------
\usepackage{xr-hyper}
\usepackage[hidelinks]{hyperref}

% External references to peer specifications
\externaldocument[reg:]{../regulations/regulations-spec}
\externaldocument[tb:]{../tb/tb-review-spec}
\externaldocument[sim:]{../simula/simula-training-spec}

% --- Fonts and listings -----------------------------------------------------
\usepackage{fontspec}
\usepackage{polyglossia}
\setmainfont{Geeza Pro}
\setmainlanguage{english}
\setotherlanguage{arabic}
\newfontfamily\arabicfont[Script=Arabic]{Mishafi}
\newfontfamily\arabicfontsf[Script=Arabic]{Mishafi}
\newfontfamily\arabicfonttt[Script=Arabic]{Mishafi}

% --- cleveref must come AFTER hyperref but keep it here (post-babel) -------
\usepackage[capitalize,noabbrev]{cleveref}

\definecolor{codebg}{gray}{0.96}
\definecolor{sapblue}{HTML}{0B5FA5}
\definecolor{sapgreen}{HTML}{00695C}
\definecolor{sapamber}{HTML}{B8360F}
\definecolor{sapgrey}{HTML}{546E7A}
\definecolor{sapmidnight}{HTML}{1C2B39}
\definecolor{sapgold}{HTML}{F0AB00}
\definecolor{apred}{HTML}{B71C1C}
\definecolor{apblue}{HTML}{1565C0}
\definecolor{apgreen}{HTML}{2E7D32}
\definecolor{apgrey}{HTML}{546E7A}
\definecolor{apamber}{HTML}{F0AB00}

\lstdefinelanguage{jsonx}{
basicstyle=\ttfamily\small,
columns=fullflexible,
breaklines=true,
showstringspaces=false,
comment=[l]{//},
morestring=[b]",
stringstyle=\color{sapgreen},
keywordstyle=\color{sapblue},
}

\onehalfspacing
\setlist{topsep=2pt,itemsep=1pt,parsep=0pt}
\setlength{\emergencystretch}{3em}
\tolerance=1000
\hbadness=10000

% --- Global Command Overrides -----------------------------------------------
\newcommand{\corpus}{}
\newcommand{\tool}{}
\newcommand{\stateid}[1]{\textsc{\lowercase{#1}}}
\newcommand{\ar}[1]{\foreignlanguage{arabic}{#1}}

% --- Glossary Helper Commands -----------------------------------------------
\definecolor{pathcolor}{HTML}{004D40}
\newcommand{\schemapath}[1]{\texttt{\textcolor{pathcolor}{#1}}}
\newcommand{\enumval}[1]{\texttt{\textcolor{sapgreen}{#1}}}
\newcommand{\fieldref}[1]{\texttt{\textcolor{sapblue}{#1}}}

\begin{document}

\title{\vspace{-2em}\Large Multi-Country Accounts Payable Processing \\
\textbf{Arabic Invoice AI Pipeline} \\
\large Implementation Specification \\
\small Ref: \texttt{AP-ARABIC-2026-v1.2}}
\author{SAP-OSS Accounts Payable Team}
\date{April 2026}

\maketitle

% --- Incorporated Executive Summary ---

\chapter*{Abstract}
\addcontentsline{toc}{chapter}{Abstract}

This specification defines the full implementation of Standard Chartered
Bank's multi-country Accounts Payable (AP) capability across the
Middle~East \& Pakistan (MEP) footprint, with particular emphasis on the
Arabic-language invoice processing pipeline sanctioned by the AI solutioning
team under the initiative titled \emph{Arabic Invoice, Peer analysis,~\& Translation}
and announced by Manoj Hemath Kumar on 2025-11-23. Three orthogonal
workstreams are specified in full:

\begin{itemize}
\item an \textbf{Arabic-invoice AI pipeline} that ingests PDF or e-mail
attachments, performs bilingual OCR and schema-guided field
extraction, applies an Arabic$\to$English translation agent,
and hands a validated record to downstream AP;
\item a \textbf{multi-country VAT compliance engine} that enforces the
regulator-specific checklists of ZATCA (Saudi Arabia), ETA
(Egypt), NBR (Bahrain), and the Qatari counterparty conventions;
\item an \textbf{AP workflow and approval engine} that formalises the
seven activities of the Bahrain \emph{Accounts Payable Departmental
Operating Instruction} (DOI) as explicit state machines.
\end{itemize}

All three workstreams consume a single canonical data schema
(\cref{ch:schema}) defined in \corpus\texttt{structured/schema/}. The
specification intentionally \emph{excludes} the downstream SAP/PeopleSoft
integration, which is covered by a separate interface document.

\paragraph{Reading order.} \Cref{ch:overview} situates the work;
\cref{ch:schema} defines the normalised record types; \cref{ch:ai} covers
the AI pipeline; \cref{ch:vat} covers the VAT engine; \cref{ch:ap}
covers the AP workflow engine; \cref{ch:controls} specifies controls,
testing, and acceptance; \cref{ch:references} is a traceability
appendix linking every requirement to a source document and a structured
record in \corpus\texttt{structured/}.

\paragraph{Conventions.} Arabic fragments are set with the
\texttt{Geeza Pro} font family, for example \ar{هيئة الزكاة والضريبة والجمارك}.
File paths are typeset in \texttt{monospace} and are relative to the
repository root. Enum values in running text appear in \textsc{small
caps}; they are defined authoritatively in
\corpus\texttt{structured/enums.yaml}.
\chapter{Overview and scope}\label{ch:overview}

\section{The problem}

Standard Chartered Bank operates Accounts Payable across a heterogeneous
regulatory and linguistic landscape. Depending on the country, an
incoming invoice may be:

\begin{itemize}
\item issued by a national e-invoicing portal (Egypt's ETA, with a
26-character \textsc{electronic\_id}; Saudi Arabia's ZATCA with
an analogous assessment reference);
\item purely Arabic (as in the two Egyptian invoices
\texttt{3CPJR5CQ\ldots K10.pdf} and \texttt{B1STE8RH\ldots SJ10.pdf}
and the Saudi VAT assessment \ar{اشعار صدور فاتورة});
\item bilingual (Qatar's QCSD / EDAA memo);
\item wholly English with pre-assigned GL coding (e.g.\ the Bahrain
DOI-covered PO flow).
\end{itemize}

Today the Arabic-language material is handled manually: a finance analyst
reads the PDF, transcribes key fields into an OTP/VAT checklist (see the
Saudi Q3 2025 checklist in \corpus\texttt{source/Saudi\ldots OTP\ldots}),
and routes it through the DOI approval chain. The November 2025 e-mail
thread from Manoj Hemath Kumar identifies this manual path as the target
for AI replacement.

\section{Goals}

The specification targets three outcomes, each with a direct lineage to
the sample corpus in \corpus\texttt{source/}:

\begin{enumerate}
\item \textbf{Structured ingestion.} Every Arabic or bilingual invoice
becomes a record conforming to
\texttt{structured/schema/invoice.schema.json}
(\cref{ch:schema}).
\item \textbf{Deterministic compliance.} Every invoice is passed
through a rule engine that emits a VAT compliance checklist
conforming to \texttt{structured/schema/vat-checklist.schema.json}
(\cref{ch:vat}), with the rule set derived from the
regulator-specific controls in the corpus.
\item \textbf{Auditable approval.} Every payment follows one of the
four state machines in \corpus\texttt{structured/workflows/} and
meets the controls enumerated in the Bahrain DOI
(\cref{ch:ap,ch:controls}).
\end{enumerate}

\section{In-scope countries and counterparties}

\begin{table}[htbp]
\centering
\caption{Countries in scope, with the controlling source document and
the structured records derived from it.}
\label{tab:countries}
\small
\begin{tabularx}{\linewidth}{@{}l l l X@{}}
\toprule
\textbf{Country} & \textbf{Authority} & \textbf{Currency} &
\textbf{Source document(s)} \\ \midrule
Egypt (EG)  & ETA   & EGP &
\texttt{3CPJR5CQ\ldots K10.pdf},\newline\texttt{B1STE8RH\ldots SJ10.pdf} \\
Saudi Arabia (SA) & ZATCA & SAR &
\ar{اشعار صدور فاتورة} (005).pdf, \newline Saudi OTP/VAT Q3'25 checklist \\
Qatar (QA)  & QTA   & QAR &
EDAA Annual Fee 2025, \newline QCSD payment instruction \\
Bahrain (BH) & NBR  & BHD &
Accounts Payable DOI \\
Iraq (IQ), UAE, PK & various & --- &
\emph{Planned}; referenced in the initiative e-mail but formally
out of scope until at least five sample invoices per jurisdiction
are added to \texttt{source/} \\ \bottomrule
\end{tabularx}
\end{table}

\section{System context}

\begin{figure}[htbp]
\centering
\begin{tikzpicture}[
node distance=9mm and 13mm,
box/.style={draw,rounded corners,minimum width=26mm,minimum height=9mm,
align=center,font=\small},
ext/.style={box,fill=blue!6},
sys/.style={box,fill=green!7},
ai/.style={box,fill=orange!10},
ap/.style={box,fill=violet!8},
arr/.style={-{Latex[length=2.2mm]},thick}
]
\node[ext] (eta)   {ETA portal\\(Egypt)};
\node[ext,right=of eta] (zatca) {ZATCA portal\\(Saudi Arabia)};
\node[ext,right=of zatca] (vendor) {Vendor PDFs\\e-mail/scan};

\node[ai,below=12mm of zatca] (ai) {Arabic-invoice\\AI pipeline\\(Ch.~\ref{ch:ai})};

\node[sys,below=of ai,xshift=-28mm] (vat) {VAT engine\\(Ch.~\ref{ch:vat})};
\node[sys,below=of ai,xshift= 28mm] (wf)  {Approval engine\\(Ch.~\ref{ch:ap})};

\node[ap,below=14mm of vat,xshift=28mm] (gl) {GL posting\\(out of scope)};

\draw[arr] (eta)    -- (ai);
\draw[arr] (zatca)  -- (ai);
\draw[arr] (vendor) -- (ai);
\draw[arr] (ai)     -- (vat);
\draw[arr] (ai)     -- (wf);
\draw[arr] (vat)    -- (gl);
\draw[arr] (wf)     -- (gl);
\end{tikzpicture}
\caption{System context. Regulator portals (blue) and vendor PDFs feed
the AI pipeline (orange), which emits canonical records to
the deterministic VAT and approval engines (green). Downstream
posting to GL is out of scope.}
\label{fig:context}
\end{figure}

\section{Out of scope}

The following are \emph{explicitly} out of scope for this document:

\begin{itemize}
\item Integration contracts with PeopleSoft PSAP/PSGL, eBBS, eOPS, and
Tax-Q. These are covered by a separate interface specification.
\item Vendor master-data on-boarding and the 48-hour screening service
itself; this document assumes a screening result is available as
a boolean precondition.
\item The encrypted template
\texttt{Business Instruction Template - V1.0 (002) (1).xlsm} and the
duplicate \texttt{EDAA\ldots[65].pdf}; neither is referenced by
any structured record.
\item Peer/competitor analysis across jurisdictions, which the e-mail
thread itself identifies as a separate initiative.
\end{itemize}

\section{Five-session engagement model}

Per the e-mail, every country instance of this system follows a
five-session agile discovery cadence
(\Cref{tab:sessions}); the AI pipeline of \cref{ch:ai} is designed to
advance with that cadence.

\begin{table}[htbp]
\centering
\caption{Five-session engagement model (codified in
\corpus\texttt{structured/enums.yaml}).}
\label{tab:sessions}
\small
\begin{tabularx}{\linewidth}{@{}l l X@{}}
\toprule
\textbf{Code} & \textbf{Name} & \textbf{Goal} \\ \midrule
\textsc{understand} & Session~1 & Map problem, define user journey, pick target. \\
\textsc{sketch}     & Session~2 & Develop detailed solution design. \\
\textsc{decide}     & Session~3 & Critique approaches, select best path. \\
\textsc{prototype}  & Session~4 & Build practical system for use case. \\
\textsc{test}       & Session~5 & Test with $\geq$5 stakeholders, gather feedback. \\
\bottomrule
\end{tabularx}
\end{table}
\chapter{Canonical data schema}\label{ch:schema}

\section{Layout of \texttt{structured/}}

The machine-readable corpus is organised as in \cref{tab:layout}.
All records validate against a JSON Schema Draft~07 file under
\corpus\texttt{structured/schema/}. The
\corpus\texttt{structured/validate.py} script walks the tree, resolves
\texttt{\$ref} entries via a local registry (so offline validation
succeeds), and exits non-zero on any failure; it is the authoritative
acceptance gate for this chapter.

\begin{table}[htbp]
\centering
\caption{Directory layout under \corpus\texttt{structured/}.}
\label{tab:layout}
\small
\begin{tabularx}{\linewidth}{@{}l X@{}}
\toprule
\textbf{Path} & \textbf{Contents} \\ \midrule
\texttt{schema/invoice.schema.json}            & Canonical invoice / assessment record. \\
\texttt{schema/vat-checklist.schema.json}      & Regulator checklist (KSA 11-point; extensible). \\
\texttt{schema/payment-instruction.schema.json}& Vendor payment memo. \\
\texttt{schema/approval-workflow.schema.json}  & Declarative state machine. \\
\texttt{schema/doi-policy.schema.json}         & DOI with activities and controls. \\
\texttt{invoices/*.json}                       & Invoice records (EG ETA ×2; SA ZATCA ×1). \\
\texttt{payment-instructions/*.json}           & Qatar QCSD/EDAA records. \\
\texttt{vat-checklists/*.json}                 & Saudi OTP/VAT Q3 2025. \\
\texttt{policies/*.json}                       & Bahrain AP DOI. \\
\texttt{workflows/*.json}                      & PO-eProc, manual-one-off, manual-recurring, VAT payment. \\
\texttt{enums.yaml}                            & Controlled vocabulary. \\
\texttt{corpus.json}                           & Flat index of every source $\leftrightarrow$ text $\leftrightarrow$ record link. \\
\texttt{validate.py}                           & Offline validator. \\ \bottomrule
\end{tabularx}
\end{table}

\section{Invoice record}

The invoice record is the central object shared by ingestion (\cref{ch:ai}),
the VAT engine (\cref{ch:vat}), and the approval engine
(\cref{ch:ap}). \Cref{tab:invoice-fields} enumerates the top-level
fields; \cref{lst:invoice-example} is a worked example taken verbatim
from \corpus\path{structured/invoices/eg-eta-3CPJR5CQ4PNNN9A4XHBB370K10.json}.

\begin{table}[htbp]
\centering
\caption{Top-level fields of \texttt{invoice.schema.json}.}
\label{tab:invoice-fields}
\small
\begin{tabularx}{\linewidth}{@{}l l c X@{}}
\toprule
\textbf{Field} & \textbf{Type} & \textbf{Req.} & \textbf{Notes} \\ \midrule
\texttt{id}               & string & \checkmark & Filename stem under \texttt{invoices/}. \\
\texttt{metadata}         & object & \checkmark & Country, authority, dates, language, status. \\
\texttt{seller}           & party  & \checkmark & Includes Arabic name, English name, \texttt{tax\_id\_type}. \\
\texttt{buyer}            & party  & \checkmark & Same shape as seller. \\
\texttt{line\_items[]}    & array  & --         & At least one of \texttt{description}, \texttt{description\_ar}, \texttt{description\_en}. \\
\texttt{totals}           & object & \checkmark & \texttt{currency}, \texttt{gross} required. \\
\texttt{tax\_period}      & object & --         & Populated for periodic assessments (ZATCA). \\
\texttt{signature}        & object & --         & Kind $\in$ \{\textsc{seal}, \textsc{wet\_ink}, \textsc{digital}, \textsc{system\_generated}\}. \\
\texttt{source\_document} & object & \checkmark & \texttt{path}, \texttt{sha256}, \texttt{ocr\_text\_path}. \\
\texttt{notes[]}          & array  & --         & Free-form analyst notes. \\
\bottomrule
\end{tabularx}
\end{table}

\begin{figure}[htbp]
\begin{lstlisting}[language=jsonx,caption={Full invoice record for the
Egypt ETA seal invoice (\texttt{eg-eta-3CPJR5CQ4PNNN9A4XHBB370K10}).
Arabic strings are UTF-8; the LaTeX engine renders them RTL.},
label=lst:invoice-example]
{
"id": "eg-eta-3CPJR5CQ4PNNN9A4XHBB370K10",
"schema_ref": "invoice.schema.json",
"metadata": {
"country": "EG",
"tax_authority": "ETA",
"document_kind": "E_INVOICE",
"electronic_id": "3CPJR5CQ4PNNN9A4XHBB370K10",
"issue_date": "2025-07-15",
"submission_date": "2025-07-15T16:10:00+00:00",
"language": "ar",
"status": "VALID"
},
"seller": {
"name_ar": "محمد شريف شرف الدين",
"tax_id": "438212282",
"tax_id_type": "ETA_TRN",
"activity_code": "1811"
},
"buyer": {
"name_ar": "بنك ستاندرد تشارترد",
"tax_id": "670307459",
"tax_id_type": "ETA_TRN"
},
"line_items": [{
"description_ar": "ختم ترودات دائرى 42 مللى سريلة",
"description_en": "Circular Trodat seal 42mm, series",
"quantity": 1, "unit_of_measure": "EA",
"unit_price": 580000.00, "line_total": 580000.00
}],
"totals": {
"currency": "EGP",
"net": 580000.00, "vat_amount": 81200.00,
"vat_rate_percent": 14, "vat_category": "FI",
"gross": 661200.00
},
"source_document": {
"path": "docs/arabic/source/3CPJR5CQ4PNNN9A4XHBB370K10.pdf",
"sha256": "2b9126...4b2c0",
"ocr_text_path": "docs/arabic/extracted-text/eg-eta-3CPJR...txt"
}
}
\end{lstlisting}
\end{figure}

Invoice \textsc{electronic\_id} for ETA must match the regex in
\corpus\texttt{structured/enums.yaml}
(\verb|^[0-9A-Z]{26}$|); the ZATCA assessment number is a decimal prefix
matching the taxpayer's \textsc{zatca\_tdn}. The \texttt{status} enum is
\{\textsc{valid}, \textsc{invalid}, \textsc{pending}, \textsc{cancelled}\};
ZATCA assessments are always \textsc{valid} on arrival.

\section{Payment instruction record}
\label{sec:payment-instruction-schema}

The payment instruction record is emitted by the approval engine
(\cref{ch:ap}) when a validated invoice and approval workflow reach a
terminal state. \Cref{tab:payment-instruction-fields} enumerates the
top-level fields of \texttt{payment-instruction.schema.json}. Two
instance records are provided in the corpus under
\texttt{structured/payment-instructions/}: the Qatar EDAA annual fee
and the QCSD payment instruction referenced in the overview chapter.

\begin{table}[htbp]
\centering
\caption{Top-level fields of \texttt{payment-instruction.schema.json}.}
\label{tab:payment-instruction-fields}
\small
\begin{tabularx}{\linewidth}{@{}l l c X@{}}
\toprule
\textbf{Field} & \textbf{Type} & \textbf{Req.} & \textbf{Notes} \\ \midrule
\texttt{id}               & string & \checkmark & Filename stem under \texttt{payment-instructions/}. \\
\texttt{metadata}         & object & \checkmark & Country, tax authority (may be null), issue/due date, status, language. \\
\texttt{payer}            & party  & \checkmark & Same party shape used in the invoice record (\cref{tab:invoice-fields}). \\
\texttt{payee}            & party  & \checkmark & Vendor receiving the funds. \\
\texttt{amount}           & object & \checkmark & \texttt{value} (number) and \texttt{currency} (ISO 4217 code). \\
\texttt{payment\_method}  & string & \checkmark & One of \texttt{bank\_transfer}, \texttt{check}, \texttt{card}, \texttt{cash}, \texttt{other}. \\
\texttt{reference}        & object & --         & Link back to the paying invoice (\texttt{invoice\_id}, \texttt{description}). \\
\texttt{approval\_chain}  & array  & --         & Links to approval-workflow records that authorised the payment. \\
\texttt{due\_date}        & string & \checkmark & ISO-8601 date. \\
\texttt{status}           & string & \checkmark & One of \texttt{draft}, \texttt{approved}, \texttt{emitted}, \texttt{settled}, \texttt{cancelled}. \\
\bottomrule
\end{tabularx}
\end{table}

\section{VAT compliance checklist}

The KSA OTP/VAT checklist is encoded as an ordered array of items
(\cref{tab:checklist-fields}). Each item records the regulator's
question, the analyst's verdict (\textsc{y}/\textsc{n}/\textsc{na}), and
the extracted value. The authoritative list is the eleven points in the
Saudi Q3 2025 checklist; NBR (Bahrain) maintains an analogous set drawn
from Activity~3 of the DOI.

\begin{table}[htbp]
\centering
\caption{Fields of a VAT checklist record.}
\label{tab:checklist-fields}
\small
\begin{tabularx}{\linewidth}{@{}l l c X@{}}
\toprule
\textbf{Field} & \textbf{Type} & \textbf{Req.} & \textbf{Notes} \\ \midrule
\texttt{period}                 & object & \checkmark & \texttt{start}, \texttt{end}, \texttt{label}. \\
\texttt{checklist[]}            & array  & \checkmark & Items indexed 1\dots$n$; each has \texttt{result} $\in$ \{Y,N,NA\}. \\
\texttt{gl\_coding}             & object & \checkmark & Chart fields (account/product/dept/OU/class). \\
\texttt{supplier}               & object & \checkmark & \texttt{supplier\_type} $\in$ \{\textsc{government}, \textsc{vendor}, \textsc{employee}, \textsc{inter\_bu}\}. \\
\texttt{tax\_type\_classification} & string & \checkmark & VAT category enum. \\
\texttt{withholding\_tax}       & object & --         & Rate and amount if applicable. \\
\texttt{psid}                   & string & --         & Preparer PSID. \\
\bottomrule
\end{tabularx}
\end{table}

\begin{figure}[htbp]
\begin{lstlisting}[language=jsonx,caption={Excerpt from
\texttt{structured/vat-checklists/sa-otp-vat-q3-2025.json}
showing items 1--5 of the KSA 11-point checklist.},label=lst:checklist-excerpt]
"checklist": [
{"n": 1, "label": "Invoice in Arabic language",                   "result": "Y"},
{"n": 2, "label": "Addressed to Standard Chartered Capital",      "result": "Y"},
{"n": 3, "label": "Date of issue",                                "result": "Y", "value": "2025-10-30"},
{"n": 4, "label": "Sequential number",                            "result": "Y", "value": "3003093648253002"},
{"n": 5, "label": "Tax Identification Number (15-digit)",         "result": "NA",
"regulation": "Supplier is ZATCA itself -- government body"}
],
\end{lstlisting}
\end{figure}

\section{Payment instruction}

Derived from the Qatar QCSD instruction. The approval chain is an array
of typed signatories; each entry carries PSID, role, and approval
timestamp --- giving a tamper-evident audit trail once hashed.
The \texttt{chart\_field} is stored as a dash-separated composite
string (e.g.\ \texttt{665011-4551-800} = account-product-operating\_unit)
because this is the convention used by the source memos.

\section{Approval workflow}

Each workflow record is a deterministic state machine. States are nodes;
transitions are edges tagged with a trigger and, where relevant, a
\texttt{required\_role} and a list of \texttt{preconditions}. The
\texttt{delegation\_of\_authority} array gates roles by LCY threshold and
is consulted by the approval engine before any transition out of a
``$x$-approved'' state. Four workflow instances are shipped in the corpus:

\begin{itemize}
\item \texttt{po-eproc} --- PO / eProcurement path (DOI Activity~3);
\item \texttt{manual-one-off} --- one-off manual payment (DOI Activity~5);
\item \texttt{manual-recurring} --- recurring manual with dispensation (DOI Activity~5);
\item \texttt{vat-payment} --- authority-counterparty (e.g.\ ZATCA) path.
\end{itemize}

\section{DOI policy}

The Bahrain Accounts Payable DOI is represented as a
\texttt{doi-policy} record whose \texttt{activities[]} array preserves
the original numbering (1--7) and whose \texttt{controls[]} array gives
each control an ID, description, pointer back to the activity, and the
risk it mitigates. This is what the approval engine consults to decide
\emph{which} workflow instance applies to a given invoice.

\section{Enumerations}

The file \corpus\texttt{structured/enums.yaml} is the single source of
truth for controlled vocabulary used across the corpus: country codes,
tax-authority codes and portal URLs, VAT category codes,
\texttt{tax\_id\_type} regexes, business-unit codes, approval-role
codes, payment methods, and the five AI-session phase codes. Any new
enum value must be added to \texttt{enums.yaml} \emph{before} being
used in a schema or record; \texttt{validate.py} enforces this.

\section{Provenance}

Each record stores the SHA-256 of its source document in
\texttt{source\_document.sha256}. Reproducing the corpus from the
originals is therefore a single command:
\begin{lstlisting}[language=bash,basicstyle=\ttfamily\small]
find docs/arabic/source -type f -print0 \
| xargs -0 shasum -a 256 \
| diff - <(jq -r '.items[] | select(.sha256) | "\(.sha256)  \(.source_path)"' \
docs/arabic/structured/corpus.json | sed 's|docs/arabic/source/||')
\end{lstlisting}
A clean diff proves the corpus is unaltered since the last
specification snapshot.

\section{Unified Schema Registry}

As of 18 April 2026, all schemas in this specification are published in the
unified schema registry at \path{docs/schema/}. Migration status is
\emph{complete} (see \path{docs/schema/registry.json} field
\texttt{migration.status}). This consolidation addresses cross-document schema
divergence identified in the specification review.

\begin{description}
\item[Registry index.] \path{docs/schema/registry.json} --- Master index of
all schemas across domains.
\item[Common types.] \path{docs/schema/common/base-types.schema.json} ---
Shared definitions (\texttt{Metadata}, \texttt{Currency},
\texttt{TaxIdentifier}, \texttt{EntityParty}) referenced via JSON
Schema \texttt{\$ref}.
\item[Consolidated enums.] \path{docs/schema/common/enums.yaml} --- Single
source of truth for enumeration values including \texttt{countries},
\texttt{vat\_category}, \texttt{supplier\_type}, and
\texttt{vat\_thresholds} used by this specification.
\item[Arabic schemas.] \path{docs/schema/arabic/} --- Domain-specific
schemas: \texttt{invoice.schema.json}, \texttt{vat-checklist.schema.json},
\texttt{vendor.schema.json}, and \texttt{ocr-result.schema.json}.
All marked \texttt{status: stable}.
\end{description}

\noindent
The per-country VAT thresholds (e.g.\ BHD~500 for Bahrain SI/FI boundary)
are defined in \path{docs/schema/common/enums.yaml\#/vat\_thresholds}
rather than in domain-specific \texttt{enums.yaml} files. Validation:

\begin{lstlisting}[language=bash]
python3 docs/schema/validate.py --registry-check
python3 docs/schema/validate.py --domain arabic \
--schema invoice.schema.json --file invoice.json
\end{lstlisting}

\noindent
CI runs \texttt{--registry-check} on every pull request touching
\path{docs/schema/} or any dependent pipeline module.
\chapter{Arabic invoice AI pipeline}\label{ch:ai}

\section{Pipeline at a glance}

\begin{figure}[htbp]
\centering
\begin{tikzpicture}[
node distance=6mm and 9mm,
stage/.style={draw,rounded corners,minimum width=32mm,minimum height=11mm,align=center,font=\small,fill=orange!10},
store/.style={draw,cylinder,shape border rotate=90,aspect=0.3,minimum width=22mm,minimum height=10mm,font=\small,fill=blue!7},
arr/.style={-{Latex[length=2.2mm]},thick}
]
\node[store] (src) {PDF / e-mail};
\node[stage,right=of src]   (s1) {1.~Intake};
\node[stage,right=of s1]    (s2) {2.~OCR};
\node[stage,below=of s2]    (s3) {3.~Field\\extraction};
\node[stage,left=of s3]     (s4) {4.~Translation};
\node[stage,left=of s4]     (s5) {5.~Validation\\\& routing};
\node[store,below=of s5]    (out) {\texttt{structured/invoices/}};

\draw[arr] (src) -- (s1);
\draw[arr] (s1)  -- (s2);
\draw[arr] (s2)  -- (s3);
\draw[arr] (s3)  -- (s4);
\draw[arr] (s4)  -- (s5);
\draw[arr] (s5)  -- (out);
\end{tikzpicture}
\caption{The five pipeline stages; each emits a well-typed payload to
the next. Only Stage~5 writes to the canonical corpus.}
\label{fig:pipeline}
\end{figure}

\section{Stage 1 -- Intake}

\paragraph{Inputs.} A source file, either:
\begin{itemize}
\item a PDF dropped into a watched folder, or
\item an e-mail attachment (e.g.\ the thread captured in
\corpus\texttt{source/FW- AI~\_ Arabic Invoice\ldots.eml}).
\end{itemize}

\paragraph{Outputs.} A new entry staged for Stage~2 and a provisional
\texttt{source\_document} object:
\begin{lstlisting}[language=jsonx]
{ "path": "docs/arabic/source/<file>", "sha256": "<sha>" }
\end{lstlisting}
The SHA-256 is computed first and compared against
\corpus\texttt{structured/corpus.json}; a match short-circuits
re-processing (idempotence) and is the duplicate-payment control
(\S\ref{ch:controls}).

\paragraph{Precondition.} The MIME detector must classify the file as
one of \texttt{application/pdf}, the OOXML Word type
(\texttt{application/vnd.openxmlformats-\allowbreak officedocument.\allowbreak wordprocessingml.\allowbreak document}),
\texttt{text/plain}, or \texttt{message/rfc822}. All other types are
rejected with reason \textsc{unsupported\_media}.

\section{Stage 2 -- OCR}

\paragraph{Command.} The corpus has been produced with:
\begin{lstlisting}[language=bash,basicstyle=\ttfamily\small]
pdftoppm -r 300 <input.pdf> work/page -png
for img in work/page-*.png; do
tesseract "$img" stdout -l ara+eng >> <output>.txt
done
\end{lstlisting}
The combined \texttt{ara+eng} language model is mandatory: invoices
routinely contain English brand names, Latin-script SKU codes, and
numerals alongside Arabic prose.

\paragraph{Quality gate.} Stage 2 emits a per-page confidence metric
(mean of tesseract's \texttt{conf} values over non-empty tokens). Pages
below \textbf{55\%} confidence are flagged for human review; below
\textbf{30\%} the file is rejected with reason \textsc{ocr\_low\_confidence}
and routed to manual processing.

\paragraph{Calibration.} The 30\,\%/55\,\% thresholds are preliminary
estimates pending production pilot validation.

\textbf{v1 status:} These values were derived from three golden-set documents,
which is insufficient for statistical confidence. The 30\,\% threshold represents
the lowest observed confidence on pages that still yielded correctable records;
55\,\% is the lowest value achieving error-free extraction. \textbf{Both thresholds
are TBD pending pilot:} production calibration requires $\geq$50 documents per
country before these values can be considered validated. The pilot will also
establish per-jurisdiction calibration as OCR quality varies with document source
(e.g., ZATCA assessments vs.\ commercial invoices).

\paragraph{Artefact.} The raw OCR text is written to
\corpus\texttt{extracted-text/<slug>.txt} and its path is recorded
in \texttt{source\_document.ocr\_text\_path}. This preserves the
training signal for future model refinement.

\section{Stage 3 -- Schema-guided field extraction}

\paragraph{Objective.} Produce a candidate record conforming to
\texttt{invoice.schema.json} from the OCR text. The extraction is
\emph{schema-guided}: the LLM prompt is constructed at run time by
serialising the schema (with \texttt{description} fields) so the
model is told exactly which JSON keys it may emit. This is the
single largest lever against hallucinated fields.

\paragraph{Prompt template.} Pseudocode; the concrete implementation
uses the Anthropic Messages API with structured-output tooling.

\begin{lstlisting}[language=bash,basicstyle=\ttfamily\small,caption={Schema-guided extraction prompt template.}]
SYSTEM: You extract tax-invoice fields into a JSON object that MUST
validate against the provided JSON Schema. If a field is not
present in the source, omit it -- NEVER invent a value.
USER:   <schema>
---
OCR TEXT (Arabic may appear in RTL order; preserve spelling):
<ocr_text>
---
Emit JSON only. Set "metadata.language" to "ar", "en", or "ar+en"
based on the dominant script.
\end{lstlisting}

\paragraph{Known-hard fields.} The following fields are the most error-prone
in the shipped corpus and must be handled with care:

\begin{table}[htbp]
\centering
\caption{Extraction pitfalls observed in the corpus and mitigations.}
\label{tab:pitfalls}
\small
\begin{tabularx}{\linewidth}{@{}l X X@{}}
\toprule
\textbf{Field} & \textbf{Pitfall} & \textbf{Mitigation} \\ \midrule
\texttt{electronic\_id}    & Tesseract maps the digit ``1'' to ``I'' in mixed alnum strings (e.g.\ \texttt{B1STE8}$\to$\texttt{BISTE8}). & Validate against the regex \verb|^[0-9A-Z]{26}$|; on mismatch, re-OCR at 600 dpi with \texttt{--psm~6}. \\
\texttt{totals.gross}      & Arabic digits (\ar{٠١٢٣٤٥٦٧٨٩}) may slip through unmapped. & Pre-tokenise: substitute Arabic-Indic digits with ASCII before parsing. \\
\texttt{vat\_rate\_percent}& Zero-rated vs exempt both show $0$. & Disambiguate with \texttt{vat\_category}: \textsc{zr} vs \textsc{ex} drawn from supplier's activity code lookup. \\
\texttt{signature.kind}    & ZATCA assessments have no signature block. & Default to \textsc{system\_generated} for \texttt{document\_kind} = \textsc{vat\_assessment}. \\
\texttt{seller.name\_en}   & Pure-Arabic invoices lack an English name. & Leave empty; Stage 4 will populate via translation. \\
\bottomrule
\end{tabularx}
\end{table}

\paragraph{Output.} A candidate record is written in memory only (never
persisted yet) and passed to Stage~4.

\section{Stage 4 -- Arabic$\to$English translation}

\paragraph{Scope.} This stage does \emph{not} translate the whole
document. It translates only the free-text fields that downstream AP
actually uses in English:

\begin{itemize}
\item \texttt{seller.name\_en} and \texttt{buyer.name\_en} if absent;
\item every \texttt{line\_items[].description\_en} if absent;
\item \texttt{seller.address\_en} and \texttt{buyer.address\_en}.
\end{itemize}

Tax IDs, amounts, dates, and enum values are \emph{never} passed through
translation; they are copied verbatim.

\paragraph{Glossary.} A fixed glossary enforces terminology consistency.
Worked examples:

\begin{table}[htbp]
\centering
\caption{Locked glossary. The right-hand column is what the translation
agent \emph{must} output.}
\label{tab:glossary}
\small
\begin{tabularx}{\linewidth}{@{}X X@{}}
\toprule
\textbf{Arabic} & \textbf{English (locked)} \\ \midrule
\ar{هيئة الزكاة والضريبة والجمارك} & Zakat, Tax and Customs Authority (ZATCA) \\
\ar{بنك ستاندرد تشارترد}           & Standard Chartered Bank \\
\ar{ضريبة القيمة المضافة}           & Value-Added Tax (VAT) \\
\ar{اشعار صدور فاتورة}              & Invoice Issuance Notification \\
\ar{فاتورة ضريبية}                  & Tax Invoice \\
\bottomrule
\end{tabularx}
\end{table}

\paragraph{Quality gate.} Back-translation is run as a sanity check:
the English output is translated back to Arabic and compared to the
source Arabic string. The metric is the \emph{normalised Levenshtein
distance} $d = \mathrm{lev}(a, b) / \max(|a|, |b|)$ computed at the
Unicode-codepoint level after Arabic NFKC normalisation (stripping
tatweel and harakat). The 0.30 ceiling was chosen because the three
golden-set Arabic strings have pairwise $d \leq 0.18$ between their
recorded Arabic and a human-verified back-translation, leaving a
$\geq 1.5\times$ safety margin. Distances above \textbf{0.30} are
flagged (\textsc{translation\_drift}) and the original Arabic is
retained as authoritative.

\section{Stage 5 -- Validation and routing}

\paragraph{Validation.} The candidate record is validated with
\texttt{validate.py} (\cref{ch:schema}). Any schema failure halts the
pipeline and the record is written to a side-queue
(\texttt{structured/invoices/rejected/}) with a \texttt{notes[]} entry
summarising the failure.

\paragraph{Routing.} On success, the pipeline consults the structured
\texttt{bh-accounts-payable-doi.json} and the \texttt{workflows/} set
to decide which state machine to instantiate:

\begin{enumerate}
\item If the counterparty is a tax authority
(\texttt{supplier\_type}~=~\textsc{government} or
\texttt{tax\_authority}~=~\textsc{zatca})
$\Rightarrow$ \texttt{vat-payment} workflow.
\item Else if there is a PO reference in \texttt{notes[]}
$\Rightarrow$ \texttt{po-eproc} workflow.
\item Else if a prior \texttt{manual-recurring} record exists for
the same \texttt{vendor\_id}
$\Rightarrow$ \texttt{manual-recurring} workflow.
\item Otherwise $\Rightarrow$ \texttt{manual-one-off} workflow.
\end{enumerate}

\paragraph{Commit.} The record is written to
\corpus\path{structured/invoices/<slug>.json} and a new item is
appended to \texttt{corpus.json}. This is the only stage in the
pipeline that mutates the corpus.

\section{Pipeline control plane}

\Cref{fig:pipeline-states} models the pipeline itself as a small state
machine; each invoice occupies exactly one of these states at any
moment. The control plane exposes three operations:
\texttt{ingest(file)}, \texttt{retry(id)}, \texttt{release(id)}.

\begin{figure}[htbp]
\centering
\begin{tikzpicture}[
node distance=6mm and 5mm,
s/.style={draw,rounded corners,minimum width=20mm,minimum height=7mm,align=center,font=\footnotesize,fill=gray!5},
t/.style={-{Latex[length=2mm]}}
]
\node[s] (new)    {\textsc{received}};
\node[s,right=of new] (ocr) {\textsc{ocr\_ok}};
\node[s,right=of ocr] (ext) {\textsc{extracted}};
\node[s,right=of ext] (xlat) {\textsc{translated}};
\node[s,right=of xlat] (val)  {\textsc{validated}};
\node[s,below=of val] (route) {\textsc{routed}};
\node[s,below=of ocr] (hitl) {\textsc{hitl\_review}};

\draw[t] (new) -- (ocr);
\draw[t] (ocr) -- (ext);
\draw[t] (ext) -- (xlat);
\draw[t] (xlat) -- (val);
\draw[t] (val) -- (route);
\draw[t] (ocr) -- node[right,font=\tiny]{conf\textless 30\%} (hitl);
\draw[t] (val) to[bend left] node[above,font=\tiny,pos=0.3]{schema fail} (hitl);
\draw[t] (hitl) -- node[below,font=\tiny,pos=0.4]{corrected} (route);
\end{tikzpicture}
\caption{Pipeline control-plane states for a single invoice.}
\label{fig:pipeline-states}
\end{figure}

\section{Evaluation harness}

\paragraph{Golden set.} The three Arabic records in
\corpus\texttt{structured/invoices/} form the regression golden set.
Every code change must re-produce byte-identical records from the
source PDFs; \texttt{validate.py} plus a \texttt{jq}-driven equality
check is the verifier.

\paragraph{Per-field metrics.} For each round of Session-5 testing the
harness reports:

\begin{itemize}
\item \textbf{Per-field exact-match accuracy} (over the 27 scalar
fields in \texttt{invoice.schema.json}).
\item \textbf{Per-field F1} on \texttt{line\_items[]} treated as sets
(descriptions matched after glossary-aware fuzzing).
\item \textbf{Human-approval rate.} Fraction of records released by
HITL with zero edits. Target $\geq$~0.75 after Session~4,
$\geq$~0.90 after Session~5.
\end{itemize}

\paragraph{Session acceptance.} The five-session model from
\cref{tab:sessions} is scheduled per country; the
\texttt{ai\_session\_phases} enum in \texttt{enums.yaml} is the system
of record. An AI rollout is complete for a country once the
Session-5 pass meets the human-approval-rate target against a cohort
of $\geq$~5 stakeholders.
\chapter{Multi-country VAT compliance engine}\label{ch:vat}

\section{Engine responsibilities}

Given a validated invoice record, the VAT engine must:

\begin{enumerate}
\item select the correct regulator and ruleset from the record's
\texttt{metadata.country} and \texttt{metadata.tax\_authority};
\item run every applicable rule, producing a
\texttt{vat-checklist.schema.json} record with one item per rule;
\item decide a final \texttt{tax\_type\_classification} (VAT category
code from \texttt{enums.yaml});
\item compute any adjustments (withholding, reverse-charge accrual);
\item refuse to release the record for payment unless every
mandatory rule emits \textsc{y} or \textsc{na}.
\end{enumerate}

Rules are declarative data, not code: the engine is a rule interpreter.
This isolates regulation change from engineering change.

\section{Rule shape}

A rule is a tuple \mbox{$(c, a, p, q, r)$} where:
\begin{itemize}
\item $c$ --- country (\textsc{sa}, \textsc{eg}, \textsc{bh}, \textsc{qa});
\item $a$ --- tax authority (\textsc{zatca}, \textsc{eta}, \textsc{nbr}, \textsc{qta});
\item $p$ --- a predicate over the invoice record (expressed as a JSON-Logic
expression or a named helper);
\item $q$ --- the checklist \emph{question} emitted when the rule fires
(human-readable string);
\item $r$ --- a \texttt{result} $\in$ \{\textsc{y}, \textsc{n}, \textsc{na}\}
computed by evaluating $p$ and optionally a secondary predicate.
\end{itemize}

The engine iterates over the ruleset for $(c,a)$ in a deterministic
order and appends each evaluated rule as an item in the output
\texttt{checklist[]}. The final \texttt{tax\_type\_classification} is
chosen by a \emph{decision function} that maps the full result vector to
a single VAT category.

\section{Saudi Arabia (ZATCA)}

\subsection{Rules}

The eleven-point ruleset is lifted verbatim from the Q3 2025 checklist
and is shown in \cref{tab:ksa-rules}. Each row is one rule; the
mapping from rule to \texttt{result} is given in the predicate column.

\begin{table}[htbp]
\centering
\caption{KSA VAT rules derived from the Q3 2025 OTP checklist.}
\label{tab:ksa-rules}
\footnotesize
\begin{tabularx}{\linewidth}{@{}c l X@{}}
\toprule
\textbf{\#} & \textbf{Question} & \textbf{Predicate / Source} \\ \midrule
1  & Invoice in Arabic language
& \texttt{metadata.language} $\in$ \{\textsc{ar}, \textsc{ar+en}\}. \\
2  & Addressed to Standard Chartered Capital
& \texttt{buyer.tax\_id} equals the entity's ZATCA TIN. \\
3  & Date of issue present
& \texttt{metadata.issue\_date} set and well-formed. \\
4  & Sequential number present
& \texttt{metadata.electronic\_id} or \texttt{metadata.invoice\_number} set. \\
5  & 15-digit TIN on supplier side
& \texttt{seller.tax\_id} matches \verb|^[0-9]{15}$|; if supplier is \textsc{government}, emit \textsc{na}. \\
6  & Nature of goods/services present
& At least one \texttt{line\_items[].description*} non-empty. \\
7  & VAT rate applied
& \texttt{totals.vat\_rate\_percent} set and $\in$ \{0, 5, 15\}. \\
8  & Net amount
& \texttt{totals.net} set. \\
9  & Gross amount
& \texttt{totals.gross} set. \\
10 & VAT amount
& \texttt{totals.vat\_amount} set and equals round($\texttt{net} \cdot r / 100$). \\
11 & Withholding tax
& \texttt{withholding\_tax.applicable}; \textsc{na} if supplier is KSA-resident. \\
\bottomrule
\end{tabularx}
\end{table}

\subsection{Decision function}

\begin{figure}[htbp]
\centering
\begin{tikzpicture}[
node distance=7mm and 11mm,
d/.style={draw,diamond,aspect=2.2,align=center,font=\footnotesize,inner sep=1pt,fill=yellow!10},
o/.style={draw,rounded corners,minimum width=20mm,minimum height=7mm,align=center,font=\small,fill=green!10},
t/.style={-{Latex[length=2mm]}}
]
\node[d] (gov)  {\textsc{government}\\counterparty?};
\node[o,right=of gov] (rc) {\textsc{rc}};
\node[d,below=of gov] (zr) {rate = 0\,\%?};
\node[o,right=of zr] (zro) {\textsc{zr}};
\node[d,below=of zr] (fi) {rate = 15\,\%?};
\node[o,right=of fi] (fio) {\textsc{fi}};
\node[o,below=of fi] (oth) {\textsc{other}};
\draw[t] (gov.east) -- node[above,font=\tiny]{yes} (rc.west);
\draw[t] (gov.south) -- node[right,font=\tiny]{no} (zr.north);
\draw[t] (zr.east) -- node[above,font=\tiny]{yes} (zro.west);
\draw[t] (zr.south) -- node[right,font=\tiny]{no} (fi.north);
\draw[t] (fi.east) -- node[above,font=\tiny]{yes} (fio.west);
\draw[t] (fi.south) -- node[right,font=\tiny]{no} (oth.north);
\end{tikzpicture}
\caption{KSA decision function. The first matching branch determines
the emitted \texttt{tax\_type\_classification}.}
\label{fig:ksa-decision}
\end{figure}

\begin{lstlisting}[language=bash,basicstyle=\ttfamily\small,caption={KSA
decision logic applied to a validated invoice record.}]
if supplier.supplier_type == GOVERNMENT
and tax_authority == ZATCA:
classification = RC         # reverse-charge / government payment
elif totals.vat_rate_percent == 0
and supplier.supplier_type != GOVERNMENT:
classification = ZR
elif totals.vat_rate_percent == 15:
classification = FI
elif totals.vat_rate_percent in {5, 10, 7, 14, 20}:
# A rate that is legal in another jurisdiction but not in KSA.
# Never silently classify: emit HOLD + human review.
classification = OTHER
verdict       = HOLD        # consumed by the AP engine
reason        = "unexpected_vat_rate_for_ksa"
else:
# Unknown rate (e.g. 3%, negative, null): always HOLD.
classification = OTHER
verdict       = HOLD
reason        = "unknown_vat_rate"
\end{lstlisting}

\subsection{Worked example}

The Saudi Q3 2025 assessment (\ar{اشعار صدور فاتورة}) has
\texttt{supplier.supplier\_type} = \textsc{government} and
\texttt{tax\_authority} = \textsc{zatca}; the decision function
correctly yields \texttt{tax\_type\_classification} = \textsc{rc}. All
eleven rules fire; rules 5 and 11 are \textsc{na} because the
counterparty is ZATCA itself and the buyer is KSA-resident. The total
\texttt{gross} = SAR~1{,}558{,}629.96 (due 2025-10-31) is posted as a
reverse-charge VAT payment through the \texttt{vat-payment} workflow.

\section{Egypt (ETA)}

\subsection{Rules}

\begin{table}[htbp]
\centering
\caption{ETA e-invoice rules.}
\label{tab:eta-rules}
\footnotesize
\begin{tabularx}{\linewidth}{@{}c l X@{}}
\toprule
\textbf{\#} & \textbf{Question} & \textbf{Predicate / Source} \\ \midrule
1 & \texttt{electronic\_id} well-formed
& Matches \verb|^[0-9A-Z]{26}$|. \\
2 & Seller TRN well-formed   & \texttt{seller.tax\_id} matches \verb|^[0-9]{9}$|. \\
3 & Buyer TRN well-formed    & \texttt{buyer.tax\_id} matches \verb|^[0-9]{9}$|. \\
4 & Activity code present    & \texttt{seller.activity\_code} set. \\
5 & Submission within 24h of issue
& \texttt{submission\_date} $-$ \texttt{issue\_date} $\leq$ 24~h. \\
6 & VAT rate consistent      & 14\% (standard) or 0\% (exempt by activity lookup). \\
7 & Line-item subtotals      & Sum of \texttt{line\_total} equals \texttt{totals.net}. \\
\bottomrule
\end{tabularx}
\end{table}

\subsection{Decision function}

\begin{lstlisting}[language=bash,basicstyle=\ttfamily\small]
if seller.activity_code in EXEMPT_ACTIVITIES:
classification = EX
elif totals.vat_rate_percent == 14:
classification = FI
elif totals.vat_rate_percent == 0:
classification = ZR
else:
classification = OTHER
\end{lstlisting}
The set \texttt{EXEMPT\_ACTIVITIES} is a configured lookup. Activity
\texttt{8890} (``Professional services'', as used by the Chapter Zero
corporate-membership invoice) is included.

\subsection{Worked examples}

\paragraph{Seal invoice (\texttt{3CPJR5CQ\ldots K10}).}
All seven rules pass; activity code \texttt{1811} is not exempt;
\texttt{totals.vat\_rate\_percent~=~14}; the decision function yields
\textsc{fi} (Full Tax Invoice). Gross \texttt{EGP~661{,}200} posts to
the PO e-Proc or manual path depending on the presence of a PO
reference.

\paragraph{Corporate membership (\texttt{B1STE8RH\ldots SJ10}).}
Activity code \texttt{8890} is in the exempt set; rule 6 still passes
because 0\% with an exempt code is consistent; decision yields \textsc{ex}.
No VAT accrues and the routing is \texttt{manual-one-off}.

\section{Bahrain (NBR)}

\subsection{Rules}

Derived from Activity 3 of the Bahrain DOI
(\corpus\path{structured/policies/bh-accounts-payable-doi.json}). The
NBR ruleset is wider than KSA's eleven points because the DOI
distinguishes ten VAT invoice types.

\begin{table}[htbp]
\centering
\caption{NBR VAT invoice types recognised in Bahrain AP.
Each row maps to a \texttt{vat\_category} enum value.}
\label{tab:nbr-types}
\footnotesize
\begin{tabularx}{\linewidth}{@{}l l X@{}}
\toprule
\textbf{DOI name} & \textbf{Enum} & \textbf{Selection criterion} \\ \midrule
Full Tax Invoice                 & \textsc{fi}  & VAT-registered supplier, $\geq$ BHD 500. \\
Simplified Tax Invoice           & \textsc{si}  & Transaction \textbf{under BHD 500}. \\
Unregistered Vendor Purchases    & \textsc{ur}  & Supplier not VAT-registered. \\
Zero Rated Purchases             & \textsc{zr}  & Transport, food, healthcare, education, construction, oil/gas. \\
Exempt Purchases                 & \textsc{ex}  & Real estate, financial services. \\
Tax Credit Notes                 & \textsc{tcn} & Negative-amount adjustment to prior invoice. \\
Reverse Charge Invoices          & \textsc{rc}  & Overseas supplier; buyer accounts for VAT. \\
GCC Tax Invoice                  & \textsc{fi}  & Intra-GCC supplier (handled as \textsc{fi} with a GCC flag). \\
Inter Business Unit Transfers    & \textsc{cbr} & Routed via GBS APU CBR team. \\
Out of Scope                     & \textsc{os}  & Pre-Jan 2019 vouchers, MPV vouchers. \\
Other                            & \textsc{other} & EWA utility bills. \\
\bottomrule
\end{tabularx}
\end{table}

\begin{figure}[htbp]
\centering
\begin{tikzpicture}[
node distance=6mm and 9mm,
d/.style={draw,diamond,aspect=2.4,align=center,font=\scriptsize,inner sep=0.5pt,fill=yellow!10},
o/.style={draw,rounded corners,minimum width=14mm,minimum height=6mm,align=center,font=\scriptsize,fill=green!10},
t/.style={-{Latex[length=1.8mm]}}
]
\node[d] (reg)  {supplier\\VAT-reg?};
\node[d,right=of reg] (os) {pre-2019?};
\node[o,right=of os] (osx) {\textsc{os}};
\node[o,below=of os] (ur) {\textsc{ur}};
\node[d,below=of reg] (ic) {inter-BU?};
\node[o,right=of ic] (cbr) {\textsc{cbr}};
\node[d,below=of ic] (neg) {credit\\note?};
\node[o,right=of neg] (tcn) {\textsc{tcn}};
\node[d,below=of neg] (rev) {overseas?};
\node[o,right=of rev] (rc) {\textsc{rc}};
\node[d,below=of rev] (zrq) {ZR\\activity?};
\node[o,right=of zrq] (zr) {\textsc{zr}};
\node[d,below=of zrq] (exq) {exempt?};
\node[o,right=of exq] (ex) {\textsc{ex}};
\node[d,below=of exq] (amt) {amount\\$<$ BHD 500?};
\node[o,right=of amt] (si) {\textsc{si}};
\node[o,below=of amt] (fi) {\textsc{fi}};

\draw[t] (reg.east) -- node[above,font=\tiny]{no} (os.west);
\draw[t] (os.east)  -- node[above,font=\tiny]{yes} (osx.west);
\draw[t] (os.south) -- node[right,font=\tiny]{no} (ur.north);
\draw[t] (reg.south) -- node[right,font=\tiny]{yes} (ic.north);
\draw[t] (ic.east) -- node[above,font=\tiny]{yes} (cbr.west);
\draw[t] (ic.south) -- node[right,font=\tiny]{no} (neg.north);
\draw[t] (neg.east) -- node[above,font=\tiny]{yes} (tcn.west);
\draw[t] (neg.south) -- node[right,font=\tiny]{no} (rev.north);
\draw[t] (rev.east) -- node[above,font=\tiny]{yes} (rc.west);
\draw[t] (rev.south) -- node[right,font=\tiny]{no} (zrq.north);
\draw[t] (zrq.east) -- node[above,font=\tiny]{yes} (zr.west);
\draw[t] (zrq.south) -- node[right,font=\tiny]{no} (exq.north);
\draw[t] (exq.east) -- node[above,font=\tiny]{yes} (ex.west);
\draw[t] (exq.south) -- node[right,font=\tiny]{no} (amt.north);
\draw[t] (amt.east) -- node[above,font=\tiny]{yes} (si.west);
\draw[t] (amt.south) -- node[right,font=\tiny]{no} (fi.north);
\end{tikzpicture}
\caption{Bahrain NBR classification cascade. The first branch
whose predicate holds emits the \texttt{vat\_category} and
short-circuits the remaining tests.}
\label{fig:nbr-decision}
\end{figure}

\subsection{BHD 500 threshold}

The DOI fixes BHD~500 as the SI/FI threshold for Bahrain. This value is
stored in the per-country engine configuration file
\texttt{vat-thresholds.yaml} (not in \texttt{enums.yaml}, which holds
only enum \emph{types}). The threshold structure is:

\begin{lstlisting}[basicstyle=\ttfamily\small]
thresholds:
BH:
si_fi_boundary: 500
currency: BHD
\end{lstlisting}

Regulator updates are applied by editing this file; a CI check validates
that all threshold values match the authoritative DOI documents.

\section{Qatar (QTA)}

Qatar does not currently impose VAT. The engine therefore runs a single
rule --- \emph{payment is to a valid counterparty with a QCSD-registered
account} --- and classifies the record as \textsc{other} unless a
future VAT regime is detected (\texttt{metadata.tax\_authority}
$\neq$ \textsc{qta}). The
\corpus\path{structured/payment-instructions/qa-qcsd-payment-instruction.json}
record is the reference.

\section{Withholding tax}

Withholding rules are modelled as a secondary ruleset keyed on
\emph{supplier residence}. The engine emits a
\texttt{withholding\_tax} object with \texttt{applicable}, \texttt{rate\_percent},
and \texttt{amount}. KSA-resident suppliers yield
\texttt{applicable~=~false} (\textsc{na}), mirroring the Q3 2025
checklist. Non-resident, non-GCC suppliers trigger a country-specific
rate (e.g.\ 5\% KSA royalty, 15\% services) that must be explicitly
configured in \texttt{enums.yaml} before the engine can emit
\textsc{y}/\textsc{n}.

\section{Engine outputs}

For every invoice the engine produces:
\begin{itemize}
\item a \texttt{vat-checklist} record stored at
\texttt{structured/vat-checklists/<slug>.json};
\item an updated \texttt{invoice} record with a populated
\texttt{totals.vat\_category};
\item a verdict of \textsc{release} or \textsc{hold} consumed by the
approval engine (\cref{ch:ap}).
\end{itemize}

A hold is produced whenever any rule resolves \textsc{n}. The engine
must never up-grade a \textsc{hold} to \textsc{release} without a human
recorded in the corresponding workflow transition.
\chapter{AP workflow and approval engine}\label{ch:ap}

\section{Engine responsibilities}

Given a validated invoice record and the VAT engine's verdict
(\cref{ch:vat}), the approval engine must:

\begin{enumerate}
\item select the correct workflow from
\corpus\texttt{structured/workflows/};
\item instantiate the state machine with the record;
\item drive transitions as approvals arrive (human or system);
\item at every transition, verify the
\texttt{delegation\_of\_authority} is respected for the current
LCY amount;
\item on reaching a terminal state, emit the payment instruction
conforming to \texttt{payment-instruction.schema.json}.
\end{enumerate}

The engine is a \emph{pure} workflow executor; it never invents fields.

\section{Activity view of the Bahrain DOI}

\Cref{tab:doi-activities} is the unabridged activity map from
\corpus\path{structured/policies/bh-accounts-payable-doi.json}; the
engine maps each activity onto one or more workflow paths.

\begin{table}[htbp]
\centering
\caption{The seven DOI activities and the workflow paths they drive.}
\label{tab:doi-activities}
\small
\begin{tabularx}{\linewidth}{@{}c l X l@{}}
\toprule
\# & \textbf{Activity} & \textbf{Description (condensed)} & \textbf{Workflow} \\ \midrule
1 & Introduction & DOI purpose \& Group alignment. & --- \\
2 & Group systems & Mandate eProc/T\&E/SODA; anything else = manual. & --- \\
3 & PO invoices  & eProc route; Tax-Q assigns POP/VAT; FCY at central-bank rate. & \texttt{po-eproc} \\
4 & Manual invoices (non-eProc) & CC + supply chain + CFO approvals; BH\_APUeInvoice mailbox. & \texttt{manual-one-off} \\
5 & In-country manual payments & One-off: MD+CFO + 48h screening. Recurring: dispensation + eOPS. & \texttt{manual-one-off}, \texttt{manual-recurring} \\
6 & CBR   & Centrally managed by GBS APU CBR team. & (central, out of this engine) \\
7 & Suspense clearing & Monthly PSGL$\leftrightarrow$eBBS reconciliation. & (batch, not per-invoice) \\
\bottomrule
\end{tabularx}
\end{table}

\section{Workflow selection}

\begin{figure}[htbp]
\centering
\begin{tikzpicture}[
node distance=8mm and 14mm,
d/.style={draw,diamond,aspect=2,align=center,font=\small,inner sep=1pt,fill=yellow!10},
w/.style={draw,rounded corners,minimum width=30mm,minimum height=8mm,align=center,font=\small,fill=violet!10},
t/.style={-{Latex[length=2mm]}}
]
\node[d] (tax)    {Tax-authority\\counterparty?};
\node[w,right=of tax] (vat) {\texttt{vat-payment}};
\node[d,below=of tax] (po)  {PO reference\\present?};
\node[w,right=of po]  (pop) {\texttt{po-eproc}};
\node[d,below=of po]  (rec) {Recurring\\vendor?};
\node[w,right=of rec] (recp){\texttt{manual-recurring}};
\node[w,below=of rec] (oneo){\texttt{manual-one-off}};

\draw[t] (tax.east) -- node[above,font=\tiny]{yes} (vat.west);
\draw[t] (tax.south) -- node[right,font=\tiny]{no} (po.north);
\draw[t] (po.east)  -- node[above,font=\tiny]{yes} (pop.west);
\draw[t] (po.south) -- node[right,font=\tiny]{no} (rec.north);
\draw[t] (rec.east) -- node[above,font=\tiny]{yes} (recp.west);
\draw[t] (rec.south)-- node[right,font=\tiny]{no} (oneo.north);
\end{tikzpicture}
\caption{Workflow selection decision tree applied at Stage~5 of the AI
pipeline (\cref{ch:ai}).}
\label{fig:workflow-select}
\end{figure}

\section{\texttt{po-eproc} state machine}

\begin{figure}[htbp]
\centering
\begin{tikzpicture}[
node distance=7mm and 6mm,
s/.style={draw,rounded corners,minimum width=22mm,minimum height=7mm,align=center,font=\scriptsize},
t/.style={-{Latex[length=2mm]}}
]
\node[s] (d)  {\textsc{draft}};
\node[s,right=of d] (po) {\textsc{po\_issued}};
\node[s,right=of po] (ir) {\textsc{invoice\_received}};
\node[s,right=of ir] (txq){\textsc{tax\_q\_review}};
\node[s,below=of txq] (posted){\textsc{posted}};
\node[s,left=of posted] (paid) {\textsc{paid}};
\node[s,below=of ir,fill=red!7] (rej) {\textsc{rejected}};

\draw[t] (d) -- node[above,font=\tiny]{CC} (po);
\draw[t] (po) -- (ir);
\draw[t] (ir) -- (txq);
\draw[t] (txq) -- node[right,font=\tiny]{VAT+POP} (posted);
\draw[t] (posted) -- node[above,font=\tiny]{APU} (paid);
\draw[t] (ir) -- node[right,font=\tiny]{reject} (rej);
\end{tikzpicture}
\caption{\texttt{po-eproc} state machine
(\corpus\texttt{structured/workflows/po-eproc.json}).}
\label{fig:po-eproc}
\end{figure}

\paragraph{Preconditions.}
\begin{itemize}
\item \texttt{tax\_q\_review}~$\to$~\texttt{posted}: VAT category
assigned and FCY revalued at central-bank rate on
\texttt{issue\_date}.
\item \texttt{posted}~$\to$~\texttt{paid}: LCY amount within DOA per
GDAM.
\end{itemize}

\section{\texttt{manual-one-off} state machine}

\Cref{fig:manual-one-off} captures Activity 5's one-off branch:
evidence of one-off nature, vendor screening refreshed within 48 hours,
MD and CFO approvals, and account coding captured in the approval
document.

\begin{figure}[htbp]
\centering
\begin{tikzpicture}[
node distance=7mm and 9mm,
s/.style={draw,rounded corners,minimum width=24mm,minimum height=7mm,align=center,font=\footnotesize},
t/.style={-{Latex[length=2mm]}}
]
\node[s] (d)  {\textsc{draft}};
\node[s,right=of d] (cc) {\textsc{cc\_approved}};
\node[s,right=of cc] (su) {\textsc{supply\_approved}};
\node[s,right=of su] (md) {\textsc{md\_approved}};
\node[s,below=of md] (cfo){\textsc{cfo\_approved}};
\node[s,left=of cfo] (scr){\textsc{screened}};
\node[s,left=of scr] (paid){\textsc{paid}};

\draw[t] (d) -- node[above,font=\tiny]{CC} (cc);
\draw[t] (cc) -- node[above,font=\tiny]{SCM} (su);
\draw[t] (su) -- node[above,font=\tiny]{MD} (md);
\draw[t] (md) -- node[right,font=\tiny]{CFO} (cfo);
\draw[t] (cfo) -- node[below,font=\tiny]{$\leq$48h} (scr);
\draw[t] (scr) -- node[below,font=\tiny]{execute} (paid);
\end{tikzpicture}
\caption{\texttt{manual-one-off}
(\corpus\texttt{structured/workflows/manual-one-off.json}).}
\label{fig:manual-one-off}
\end{figure}

\section{\texttt{manual-recurring} state machine}

\paragraph{Precondition.} A written dispensation from the policy owner
is \emph{mandatory} before entry. This is enforced by the engine: the
record cannot leave \textsc{draft} until a
\texttt{dispensation\_document\_path} is recorded in \texttt{notes[]}.
Post-payment, the record must reach \textsc{quarterly\_reviewed} within
the following quarter, triggered by the CFO's quarterly log retrieval.

\section{\texttt{vat-payment} state machine}

Used when the counterparty is a tax authority
(\texttt{supplier.supplier\_type} = \textsc{government}). The KSA Q3
2025 assessment walks this workflow from
\textsc{assessment\_received} $\to$ \textsc{checklist\_validated}
(upon 11/11 rules resolved) $\to$ \textsc{cfo\_approved} $\to$
\textsc{paid} $\to$ \textsc{posted}.

\begin{figure}[htbp]
\centering
\begin{tikzpicture}[
node distance=5mm and 5mm,
s/.style={draw,rounded corners,minimum width=22mm,minimum height=7mm,align=center,font=\scriptsize,fill=teal!8},
t/.style={-{Latex[length=2mm]}}
]
\node[s] (a)  {\textsc{assessment\_received}};
\node[s,right=of a] (cv) {\textsc{checklist\_validated}};
\node[s,right=of cv] (ca) {\textsc{cfo\_approved}};
\node[s,below=of ca] (p)  {\textsc{paid}};
\node[s,left=of p]   (po) {\textsc{posted}};

\draw[t] (a) -- node[above,font=\tiny]{11/11 rules} (cv);
\draw[t] (cv) -- node[above,font=\tiny]{CFO} (ca);
\draw[t] (ca) -- node[right,font=\tiny]{SAMA rail} (p);
\draw[t] (p) -- node[above,font=\tiny]{GL post} (po);
\end{tikzpicture}
\caption{\texttt{vat-payment}: linear path walked by the ZATCA Q3
2025 assessment. Entry requires every VAT checklist rule to resolve
to \textsc{y} or \textsc{na}.}
\label{fig:vat-payment}
\end{figure}

\section{Delegation-of-authority}

The engine enforces DOA as a linear scan of the workflow's
\texttt{delegation\_of\_authority} array, sorted ascending by
\texttt{max\_amount}. The role required is the first entry whose
\texttt{max\_amount} is $\geq$ the invoice's LCY value. An FX
conversion (at the central-bank rate on \texttt{issue\_date}) is
applied before comparison.

\Cref{tab:doa} shows the shipped thresholds.

\begin{table}[htbp]
\centering
\caption{Default DOA thresholds per workflow.}
\label{tab:doa}
\small
\begin{tabularx}{\linewidth}{@{}l l r l@{}}
\toprule
\textbf{Workflow} & \textbf{Currency} & \textbf{Max LCY} & \textbf{Required role} \\ \midrule
\texttt{po-eproc}        & BHD &    10{,}000 & CC\_OWNER \\
\texttt{po-eproc}        & BHD &   100{,}000 & SUPPLY\_CHAIN\_MGR \\
\texttt{po-eproc}        & BHD & 1{,}000{,}000 & CFO \\
\texttt{manual-one-off}  & BHD &      500 & SUPPLY\_CHAIN\_MGR \\
\texttt{manual-one-off}  & BHD &   50{,}000 & CFO \\
\texttt{manual-one-off}  & BHD & 1{,}000{,}000 & MD \\
\texttt{manual-recurring}& BHD & 1{,}000{,}000 & CFO \\
\texttt{vat-payment}     & SAR & 10{,}000{,}000 & CFO \\
\bottomrule
\end{tabularx}
\end{table}

\section{Exceptions}

The shipped workflows carry an \texttt{exceptions[]} array. The most
important entry is the GSAM legal/professional fees exception
(\texttt{manual-one-off}): these are \emph{not} covered by the default
GDAM approval limits and must be routed to a GSAM-specific approver. The
engine consults \texttt{exceptions[]} before applying DOA.

\section{Payment-suspense reconciliation}

Activity 7 is operated as a \emph{batch} procedure; the approval engine
does not execute it per invoice. Reconciliation compares the PSGL
extract against the eBBS extract on the accounts listed in the DOI:

\begin{itemize}
\item DBU BHD suspense: \texttt{09-0906743-50}
\item DBU USD suspense: \texttt{09-0906743-50}
\end{itemize}

The aging escalation grid (30/60/90 day buckets) drives automated
e-mails to the DBU finance manager; items aged $>$90~days escalate to
the CFO.

\section{Auditability}

Every transition records a structured event:
\texttt{(workflow\_id, from, to, actor\_psid, timestamp, evidence\_refs)}.
The concatenation of events is a tamper-evident log because each event
hashes its predecessor. This is what downstream auditors consume; it
is \emph{not} part of the record schema but sits alongside the
\texttt{structured/} corpus in an append-only \texttt{events/} stream.
The event contract specified here is in scope (\cref{ch:controls},
``Audit event log''); the storage substrate, retention policy, and
restoration runbook are owned by GBS APU Audit and documented in
\emph{SCB-APU-AUDIT-001}.
\chapter{Controls, testing, and acceptance}\label{ch:controls}

\section{Control catalogue}

The Bahrain DOI's seven controls
(\corpus\path{structured/policies/bh-accounts-payable-doi.json},
\texttt{controls[]}) are reproduced in \cref{tab:controls} with the
engine mechanism that enforces each.

\begin{table}[htbp]
\centering
\caption{Control catalogue and enforcement mechanisms.}
\label{tab:controls}
\small
\begin{tabularx}{\linewidth}{@{}l l X X@{}}
\toprule
\textbf{ID} & \textbf{Activity} & \textbf{Description} & \textbf{Enforced by} \\ \midrule
C1 & 4 & All required approvals present before payment release. & Workflow engine cannot transition to \textsc{paid} until every required-role approval has fired. \\
C2 & 5 & Vendor screening revalidated within 48~hours of payment. & Precondition on \texttt{cfo\_approved}$\to$\textsc{screened} transition (\texttt{manual-one-off}). \\
C3 & 4 & Amount payable within DOA per GDAM. & Linear DOA scan before any approval transition. \\
C4 & 5 & Dispensation in place for recurring manual payments. & Precondition on \textsc{draft}$\to$\textsc{dispensation\_ok} in \texttt{manual-recurring}. \\
C5 & 7 & Monthly PSGL$\leftrightarrow$eBBS reconciliation cleared within aging thresholds. & Batch job; escalation grid e-mails DBU FM / CFO. \\
C6 & 3 & FCY invoices revalued at central-bank rate on invoice date. & VAT engine runs FCY$\to$LCY \emph{before} DOA check. \\
C7 & 4 & Duplicate-payment check. & Intake Stage~1 short-circuits on matching \texttt{source\_document.sha256}. \\
\bottomrule
\end{tabularx}
\end{table}

\begin{figure}[htbp]
\centering
\begin{tikzpicture}[
node distance=5mm and 6mm,
stage/.style={draw,rounded corners,minimum width=22mm,minimum height=8mm,align=center,font=\scriptsize,fill=orange!10},
ctrl/.style={draw,ellipse,minimum width=10mm,minimum height=6mm,font=\tiny,fill=red!10},
t/.style={-{Latex[length=1.8mm]}},
c/.style={-{Latex[length=1.6mm]},dashed,red!60!black}
]
\node[stage] (intake) {Intake};
\node[stage,right=of intake]  (ocr)  {OCR};
\node[stage,right=of ocr]     (ext)  {Extract};
\node[stage,right=of ext]     (vat)  {VAT engine};
\node[stage,below=of vat]     (doa)  {DOA scan};
\node[stage,left=of doa]      (apr)  {Approval};
\node[stage,left=of apr]      (pay)  {Pay};
\node[stage,left=of pay]      (rec)  {Reconcile};

\draw[t] (intake) -- (ocr);
\draw[t] (ocr) -- (ext);
\draw[t] (ext) -- (vat);
\draw[t] (vat) -- (doa);
\draw[t] (doa) -- (apr);
\draw[t] (apr) -- (pay);
\draw[t] (pay) -- (rec);

\node[ctrl,above=2.5mm of intake] (c7)  {C7};
\node[ctrl,above=2.5mm of vat]    (c6)  {C6};
\node[ctrl,above=2.5mm of doa]    (c3)  {C3};
\node[ctrl,below=2.5mm of apr]    (c1)  {C1};
\node[ctrl,below=2.5mm of pay]    (c2)  {C2};
\node[ctrl,below=2.5mm of pay,xshift=-14mm] (c4) {C4};
\node[ctrl,below=2.5mm of rec]    (c5)  {C5};

\draw[c] (c7) -- (intake);
\draw[c] (c6) -- (vat);
\draw[c] (c3) -- (doa);
\draw[c] (c1) -- (apr);
\draw[c] (c2) -- (pay);
\draw[c] (c4) -- (pay);
\draw[c] (c5) -- (rec);
\end{tikzpicture}
\caption{Control enforcement points along the end-to-end invoice
lifecycle. Dashed arrows show where each DOI control
(\cref{tab:controls}) is evaluated.}
\label{fig:controls-map}
\end{figure}

\section{Testing strategy}

Testing follows the five-session engagement model
(\cref{tab:sessions}). For each country rollout:

\begin{description}
\item[Session 1 --- Understand.] Collect $\geq$5 sample invoices
(regulator-issued if available), normalise into the corpus,
confirm per-country ruleset is complete.
\item[Session 2 --- Sketch.] Walk the end-to-end pipeline on paper
using one sample invoice. Identify OCR pitfalls
(\cref{tab:pitfalls}) specific to the country.
\item[Session 3 --- Decide.] Freeze the rule precedence and the
decision function for the VAT engine
(\S\ref{ch:vat}).
\item[Session 4 --- Prototype.] Run the full pipeline on the cohort;
review each extracted record with the local finance team.
Human-approval rate target $\geq$~0.75.
\item[Session 5 --- Testing.] Independent run with $\geq$~5
stakeholders; target $\geq$~0.90 human-approval with
zero-edit rate measured.
\end{description}

\section{Regression golden set}

The three Arabic records already in
\corpus\texttt{structured/invoices/} (Egypt ETA $\times$ 2, Saudi
ZATCA $\times$ 1) together with the KSA checklist and the two Qatar
payment instructions constitute the regression golden set.

\paragraph{Invariants.} Every change to the pipeline must preserve:
\begin{enumerate}
\item byte-identical \texttt{source\_document.sha256} for every
corpus record (source files are immutable);
\item 100\,\% schema validity under \texttt{validate.py};
\item identical \texttt{tax\_type\_classification} for every
invoice record;
\item identical selected workflow for every record.
\end{enumerate}

\paragraph{Runner.} \texttt{validate.py} plus the following
\texttt{jq} equality check is the regression harness:

\begin{lstlisting}[language=bash,basicstyle=\ttfamily\small]
# Per-field equality of every golden-set record between a reference
# snapshot and the pipeline output.
diff \
<(jq -S . docs/arabic/structured/invoices/*.json) \
<(for p in docs/arabic/source/*.pdf; do
./pipeline extract "$p"
done | jq -S .)
\end{lstlisting}

A non-empty diff is a regression.

\section{Acceptance criteria per workstream}

\paragraph{AI pipeline (\cref{ch:ai}).}
\begin{itemize}
\item Corpus integrity check passes (every \texttt{corpus.json} path
exists; every SHA-256 matches).
\item For every source PDF: OCR confidence $\geq$~30\% on every page,
schema-valid record emitted.
\item Translation back-translation distance $\leq$~0.30 on every
non-empty \texttt{description\_en}.
\end{itemize}

\paragraph{VAT engine (\cref{ch:vat}).}
\begin{itemize}
\item KSA: all 11 checklist rules evaluated; decision function
emits expected VAT category for every golden invoice.
\item Egypt: \texttt{EXEMPT\_ACTIVITIES} lookup table complete for
activities present in the corpus.
\item Bahrain: each of the 11 NBR invoice types has a corresponding
enum code.
\end{itemize}

\paragraph{AP engine (\cref{ch:ap}).}
\begin{itemize}
\item Every workflow record loadable and schema-valid.
\item DOA scan deterministic under round-trip FX conversion.
\item Workflow-selection decision tree (\cref{fig:workflow-select})
yields the documented workflow for every corpus record.
\end{itemize}

\section{Traceability matrix}

\Cref{tab:trace} links each specification section to the source
document that justifies it and to the structured record that realises
it. Every row must be present for a full audit pass.

\begin{table}[htbp]
\centering
\caption{Traceability of specification to source document and
structured record.}
\label{tab:trace}
\footnotesize
\begin{tabularx}{\linewidth}{@{}l X X X@{}}
\toprule
\textbf{Section} & \textbf{Requirement} & \textbf{Source} & \textbf{Structured record} \\ \midrule
\cref{ch:overview} & Countries in scope & Manoj e-mail, 2025-11-23 & \path{enums.yaml} (countries) \\
\cref{ch:schema}   & Canonical invoice fields & EG ETA PDFs, ZATCA PDF & \path{schema/invoice.schema.json} \\
\cref{ch:schema}   & 11-point checklist shape & Saudi OTP/VAT Q3 2025 & \path{schema/vat-checklist.schema.json} \\
\cref{ch:ai}       & OCR command line & --- (new) & \path{extracted-text/*.txt} \\
\cref{ch:ai}       & Translation glossary & \ar{اشعار صدور فاتورة} (005).pdf & \path{invoices/sa-zatca-ishaar-005.json} \\
\cref{ch:vat}      & KSA rules (1--11) & Saudi OTP/VAT Q3 2025 & \path{vat-checklists/sa-otp-vat-q3-2025.json} \\
\cref{ch:vat}      & ETA 26-char ID regex & EG ETA PDFs & \path{enums.yaml} (tax\_authorities.ETA) \\
\cref{ch:vat}      & NBR 11 invoice types & AP DOI Activity 3 & \path{policies/bh-accounts-payable-doi.json} \\
\cref{ch:ap}       & Activities 1--7 & AP DOI & \path{policies/bh-accounts-payable-doi.json} \\
\cref{ch:ap}       & \texttt{po-eproc} state machine & AP DOI Activity 3 & \path{workflows/po-eproc.json} \\
\cref{ch:ap}       & \texttt{manual-one-off} state machine & AP DOI Activity 5 (one-off) & \path{workflows/manual-one-off.json} \\
\cref{ch:ap}       & \texttt{manual-recurring} state machine & AP DOI Activity 5 (recurring) & \path{workflows/manual-recurring.json} \\
\cref{ch:ap}       & \texttt{vat-payment} state machine & KSA Q3 2025 flow + ZATCA assessment & \path{workflows/vat-payment.json} \\
\cref{ch:controls} & Controls C1--C7 & AP DOI + KSA checklist + observations & \path{policies/bh-accounts-payable-doi.json} \\
\bottomrule
\end{tabularx}
\end{table}

\section{Operational readiness}

A country is declared \emph{production-ready} once:
\begin{enumerate}
\item the regression golden set runs clean for 30 consecutive days
in the \textsc{uat} environment;
\item Session~5 acceptance meets the human-approval-rate target
(\cref{ch:ai});
\item every VAT rule for the country is present in the ruleset and
has at least one passing test case;
\item the country CFO signs off on the workflow configuration,
recorded as a \texttt{CFO} role entry on the \texttt{vat-payment}
or \texttt{manual-one-off} workflow golden record.
\end{enumerate}

\section{Deployment and rollback}

\paragraph{Environments.} Three are in scope:
\begin{description}
\item[\textsc{dev}.] Developer workstations and CI; reads the corpus
from the Git working tree. No live invoices.
\item[\textsc{uat}.] Shared test environment reading a frozen copy of
the corpus; live DOI activity stops at \textsc{cfo\_approved}
(no SAMA/SWIFT rail). Country finance and GBS APU run Session
4/5 here.
\item[\textsc{prod}.] Live environment; the only one authorised to
transition to \textsc{paid}.
\end{description}

\paragraph{Sign-off matrix.} Each promotion requires a two-party
approval recorded as an event in the \texttt{events/} stream
(\cref{ch:ap}):

\begin{tabularx}{\linewidth}{@{}l l X@{}}
\toprule
\textbf{Transition} & \textbf{Approvers} & \textbf{Evidence} \\ \midrule
\textsc{dev}$\to$\textsc{uat}  & Tech lead + QA lead           & Green CI build, clean regression. \\
\textsc{uat}$\to$\textsc{prod} & Country CFO + GBS APU head    & Session 5 acceptance report; 30-day UAT stability; CFO workflow golden record. \\
\bottomrule
\end{tabularx}

\paragraph{Rollback.} Deployments are blue/green. The engine, rules,
and schemas are versioned with a content-addressed tag (SHA-256 of
\texttt{structured/corpus.json} at release time). Rollback is a
one-command revert to the previous tag; no database migration is
required because the canonical corpus is append-only. Any in-flight
invoice whose workflow record references a schema version that was
rolled back is automatically held in \textsc{awaiting\_review}.

\paragraph{Change control.} Ruleset edits (\cref{ch:vat}) and
workflow edits (\cref{ch:ap}) are routed through the same two-party
sign-off; schema edits require an additional sign-off from the AI
pipeline owner because they invalidate the regression golden set
(\S\ref{ch:ai}).

\section{Audit event log}

The approval engine emits an append-only event per transition
(\texttt{workflow\_id}, \texttt{from}, \texttt{to},
\texttt{actor\_psid}, \texttt{timestamp}, \texttt{evidence\_refs}),
hashing its predecessor so the stream is tamper-evident. The stream
is \emph{in scope} for compliance purposes but the storage substrate
(append-only object store, retention policy, restoration runbook) is
owned by the GBS APU audit team and specified in
\emph{SCB-APU-AUDIT-001}, referenced but not reproduced here. The
engine guarantees:

\begin{itemize}
\item every workflow transition in \cref{ch:ap} emits exactly one
event;
\item events are linearly hash-chained;
\item the event stream is the sole source of truth for DOA
compliance (controls C1--C4 in \cref{tab:controls}).
\end{itemize}
\chapter{References}\label{ch:references}

This appendix lists every file that backs a normative statement in the
specification. Paths are relative to the repository root. SHA-256 hashes
are taken from \corpus\texttt{structured/corpus.json} and are the
ground truth for the provenance check of \cref{ch:schema}.

\section{Source documents}

\footnotesize
\begin{longtable}{@{}p{3.4cm}p{8cm}p{3cm}@{}}
\toprule
\textbf{Short name} & \textbf{Path} & \textbf{SHA-256 (prefix)} \\ \midrule
\endfirsthead
\toprule
\textbf{Short name} & \textbf{Path} & \textbf{SHA-256 (prefix)} \\ \midrule
\endhead
EG ETA seal invoice & \path{docs/arabic/source/3CPJR5CQ4PNNN9A4XHBB370K10.pdf} & \texttt{2b912687fa94\ldots} \\
EG ETA corp membership & \path{docs/arabic/source/B1STE8RHNWZY9TYEXDS6WSZJ10.pdf} & \texttt{18eb35671131\ldots} \\
SA ZATCA assessment & \path{docs/arabic/source/} \ar{اشعار صدور فاتورة} (005).pdf & \texttt{cbbb3f702ad5\ldots} \\
QA EDAA annual fee & \path{docs/arabic/source/EDAA Annual Fee Invoice 2025.pdf} & \texttt{cf91c2a43574\ldots} \\
QA QCSD payment instr. & \path{docs/arabic/source/Instruction for payment to vendor (QCSD).pdf} & \texttt{39ed62690d0f\ldots} \\
SA OTP VAT Q3 2025 & \path{docs/arabic/source/Saudi - OTP Entry processing- VAT Q3'25 (1).docx} & \texttt{3a6c8f37528f\ldots} \\
BH AP DOI & \path{docs/arabic/source/Accounts Payable DOI.docx} & \texttt{128ff6792798\ldots} \\
AI initiative e-mail & \path{docs/arabic/source/FW- AI _ Arabic Invoice...eml} & \texttt{e01ba78d5ca5\ldots} \\
\bottomrule
\end{longtable}
\normalsize

\section{Derived text}

Every Arabic PDF has a UTF-8 OCR artefact in
\corpus\texttt{extracted-text/}, produced with
\texttt{pdftoppm -r 300 | tesseract -l ara+eng}. English PDFs use
\texttt{pdftotext -layout}; Word documents use
\texttt{pandoc -t plain}; the e-mail is decomposed by a small Python
script driven by the standard-library \texttt{email} module.

\section{Structured records}

All of the following are JSON records validated by
\corpus\texttt{structured/validate.py}.

\footnotesize
\begin{longtable}{@{}p{4cm}p{10cm}@{}}
\toprule
\textbf{Type} & \textbf{Path} \\ \midrule
\endhead
Invoice & \path{structured/invoices/eg-eta-3CPJR5CQ4PNNN9A4XHBB370K10.json} \\
Invoice & \path{structured/invoices/eg-eta-B1STE8RHNWZY9TYEXDS6WSZJ10.json} \\
Invoice (assessment) & \path{structured/invoices/sa-zatca-ishaar-005.json} \\
Payment instruction & \path{structured/payment-instructions/qa-edaa-annual-fee-2025.json} \\
Payment instruction & \path{structured/payment-instructions/qa-qcsd-payment-instruction.json} \\
VAT checklist & \path{structured/vat-checklists/sa-otp-vat-q3-2025.json} \\
DOI policy & \path{structured/policies/bh-accounts-payable-doi.json} \\
Workflow & \path{structured/workflows/po-eproc.json} \\
Workflow & \path{structured/workflows/manual-one-off.json} \\
Workflow & \path{structured/workflows/manual-recurring.json} \\
Workflow & \path{structured/workflows/vat-payment.json} \\
Index & \path{structured/corpus.json} \\
Enumerations & \path{structured/enums.yaml} \\
\bottomrule
\end{longtable}
\normalsize

\section{Tools}

\begin{tabularx}{\linewidth}{@{}l X@{}}
\toprule
\textbf{Tool} & \textbf{Purpose} \\ \midrule
\texttt{pdftoppm}    & Rasterise PDF pages to 300~dpi PNG for OCR. \\
\texttt{tesseract}   & Bilingual OCR with \texttt{-l ara+eng}. \\
\texttt{pdftotext}   & Direct text extraction for born-digital English PDFs. \\
\texttt{pandoc}      & Convert \texttt{.docx} to plain text. \\
Python \texttt{email}& Parse the \texttt{.eml} envelope and MIME parts. \\
\texttt{jq}          & Query and cross-check \texttt{corpus.json}. \\
\texttt{jsonschema} / \texttt{referencing} (Python) & Offline schema validation with a local registry. \\
\texttt{xelatex} / \texttt{tectonic} & Build this document (polyglossia, bidi support). \\
\bottomrule
\end{tabularx}

\section{Minimal corpus excerpt}

The entire corpus is in-tree under \corpus; for readers without
access to the Git repository, \cref{lst:corpus-excerpt} reproduces
the first three items of \corpus\texttt{structured/corpus.json}
verbatim. This is a sufficient starting point to reconstruct the
shape of the index; the full file contains 12 items covering every
source in \cref{tab:layout}.

\begin{figure}[htbp]
\begin{lstlisting}[language=jsonx,caption={First three items of
\texttt{corpus.json}; identical SHA-256 values are reproduced in
\cref{ch:references}.},label=lst:corpus-excerpt]
{
"schema_version": "1.0",
"items": [
{
"id": "eg-eta-3CPJR5CQ4PNNN9A4XHBB370K10",
"type": "invoice",
"country": "EG",
"tax_authority": "ETA",
"sha256": "2b912687fa94...",
"source_path":     "docs/arabic/source/3CPJR5CQ...K10.pdf",
"text_path":       "docs/arabic/extracted-text/eg-eta-...K10.txt",
"structured_path": "docs/arabic/structured/invoices/eg-eta-...K10.json"
},
{
"id": "sa-zatca-ishaar-005",
"type": "invoice",
"country": "SA",
"tax_authority": "ZATCA",
"sha256": "cbbb3f702ad5...",
"source_path":     "docs/arabic/source/<arabic-name>.pdf",
"text_path":       "docs/arabic/extracted-text/sa-zatca-ishaar-005.txt",
"structured_path": "docs/arabic/structured/invoices/sa-zatca-ishaar-005.json"
},
{
"id": "qa-edaa-annual-fee-2025",
"type": "payment_instruction",
"country": "QA",
"tax_authority": null,
"sha256": "cf91c2a43574..."
}
]
}
\end{lstlisting}
\end{figure}

\section{Out-of-scope assets}

The following files sit in
\corpus\texttt{source/} but are \emph{not} part of the corpus and
are not referenced by any record:

\begin{itemize}
\item \texttt{Business Instruction Template - V1.0 (002) (1).xlsm} ---
encrypted CDFV2 workbook; cannot be opened without the owner's
password.
\item \texttt{EDAA Annual Fee Invoice 2025[65].pdf} --- byte-identical
duplicate of \texttt{EDAA Annual Fee Invoice 2025.pdf}.
\end{itemize}
\chapter{VAT Module Implementation Specification}\label{ch:vat-impl}

This chapter provides the technical implementation specification for the VAT
processing module, addressing the review finding that ``VAT engine rule format
is underspecified.'' The specification enables developers to build conformant
implementations.

\section{Module architecture}

The VAT module comprises four components:

\begin{enumerate}
\item \textbf{RulesEngine} --- evaluates JSON-Logic expressions against invoice
data and produces checklist items;
\item \textbf{RegistrationValidator} --- validates VAT registration numbers
against country-specific formats and (optionally) authority APIs;
\item \textbf{Calculator} --- computes expected VAT amounts given rates and
subtotals;
\item \textbf{ActivityLookup} --- maps business activity codes to VAT
treatment categories.
\end{enumerate}

\begin{figure}[htbp]
\centering
\begin{tikzpicture}[
node distance=8mm and 15mm,
box/.style={draw,rounded corners,minimum width=28mm,minimum height=10mm,align=center,font=\small,fill=blue!5},
db/.style={draw,cylinder,shape border rotate=90,aspect=0.25,minimum width=18mm,minimum height=12mm,fill=gray!10,font=\small},
t/.style={-{Latex[length=2mm]}}
]
\node[box] (eng) {RulesEngine};
\node[box,below left=of eng,align=center] (reg) {Registration\\Validator};
\node[box,below=of eng] (calc) {Calculator};
\node[box,below right=of eng,align=center] (act) {Activity\\Lookup};
\node[db,above=of eng] (rules) {rules.json};
\node[db,above left=of eng,xshift=-5mm,align=center] (inv) {invoice\\record};
\node[db,above right=of eng,xshift=5mm,align=center] (out) {checklist\\record};

\draw[t] (inv.south) -- (eng.north west);
\draw[t] (rules.south) -- (eng.north);
\draw[t] (eng.north east) -- (out.south);
\draw[t,dashed] (eng.south west) -- (reg.north east);
\draw[t,dashed] (eng.south) -- (calc.north);
\draw[t,dashed] (eng.south east) -- (act.north west);
\end{tikzpicture}
\caption{VAT module component diagram. Dashed arrows indicate helper
invocations from the rules engine.}
\label{fig:vat-arch}
\end{figure}

\section{API contracts}

Each helper function has a defined API contract. Implementations \textbf{must}
conform to these signatures to be invoked by the rules engine.

\subsection{Registration validation}

\begin{lstlisting}[language=Python,basicstyle=\ttfamily\small,caption={Registration validator API contract.}]
def check_vat_registration_validity(
tax_id: str,
country: str  # ISO 3166-1 alpha-2
) -> ValidationResult:
"""
Validate a VAT registration number.

Returns:
ValidationResult with fields:
- valid: bool
- format_valid: bool
- authority_confirmed: Optional[bool]
- error_code: Optional[str]
"""
\end{lstlisting}

\begin{table}[htbp]
\centering
\caption{VAT registration number formats by country.}
\label{tab:vat-formats}
\footnotesize
\begin{tabularx}{\linewidth}{@{}l l X@{}}
\toprule
\textbf{Country} & \textbf{Regex} & \textbf{Notes} \\ \midrule
SA & \verb|^[0-9]{15}$| & 15-digit ZATCA TIN \\
EG & \verb|^[0-9]{9}$| & 9-digit ETA TRN \\
BH & \verb|^[0-9]{9,15}$| & NBR VAT number (variable length) \\
QA & \verb|^[A-Z0-9]{10}$| & QCSD registration (non-VAT) \\
\bottomrule
\end{tabularx}
\end{table}

\subsection{VAT calculation}

\begin{lstlisting}[language=Python,basicstyle=\ttfamily\small,caption={VAT calculator API contract.}]
def calculate_vat_amount(
subtotal: Decimal,
vat_rate_percent: Decimal
) -> CalculationResult:
"""
Calculate expected VAT amount.

Returns:
CalculationResult with fields:
- vat_amount: Decimal (rounded to 2 decimal places)
- gross_amount: Decimal
- calculation_method: str  # "standard" | "reverse"
"""
\end{lstlisting}

The calculator \textbf{must} use banker's rounding (ROUND\_HALF\_EVEN) and
\textbf{must} handle the following edge cases:

\begin{enumerate}
\item Negative subtotals (credit notes): return negative VAT amount
\item Zero rate: return VAT amount of exactly 0.00
\item Rate $>$ 100\%: reject with error code \texttt{INVALID\_RATE}
\end{enumerate}

\subsection{Activity lookup}

\begin{lstlisting}[language=Python,basicstyle=\ttfamily\small,caption={Activity lookup API contract.}]
def lookup_activity_code_treatment(
activity_code: str,
country: str  # ISO 3166-1 alpha-2
) -> TreatmentResult:
"""
Determine VAT treatment from business activity code.

Returns:
TreatmentResult with fields:
- vat_category: str  # ZR, EX, FI, etc.
- rate_percent: Optional[Decimal]
- source: str  # regulation reference
"""
\end{lstlisting}

\begin{table}[htbp]
\centering
\caption{Activity code to VAT treatment mapping (Egypt, partial).}
\label{tab:activity-lookup}
\footnotesize
\begin{tabularx}{\linewidth}{@{}l l l X@{}}
\toprule
\textbf{Code} & \textbf{Description} & \textbf{Treatment} & \textbf{Source} \\ \midrule
1811 & Printing services & FI (14\%) & ETA Circular 2019/1 \\
8890 & Professional services & EX (exempt) & ETA Circular 2020/3 \\
4911 & Passenger rail transport & ZR (0\%) & VAT Law Art. 23 \\
6411 & Banking services & EX (exempt) & VAT Law Art. 29 \\
\bottomrule
\end{tabularx}
\end{table}

\section{JSON-Logic expression specification}

The rules engine evaluates predicates expressed in JSON-Logic format
(\url{https://jsonlogic.com}). This section specifies the supported
operators and variable paths.

\subsection{Supported operators}

\begin{table}[htbp]
\centering
\caption{JSON-Logic operators supported by the rules engine.}
\label{tab:jsonlogic-ops}
\footnotesize
\begin{tabularx}{\linewidth}{@{}l l X@{}}
\toprule
\textbf{Category} & \textbf{Operators} & \textbf{Notes} \\ \midrule
Comparison & \texttt{==}, \texttt{===}, \texttt{!=}, \texttt{!==} & Strict equality recommended \\
& \texttt{<}, \texttt{<=}, \texttt{>}, \texttt{>=} & Numeric comparison \\
Logic      & \texttt{and}, \texttt{or}, \texttt{!}, \texttt{!!} & Boolean operations \\
& \texttt{if} & Ternary conditional \\
String     & \texttt{in}, \texttt{cat}, \texttt{substr} & String operations \\
Array      & \texttt{in}, \texttt{merge}, \texttt{map}, \texttt{filter} & Array operations \\
Arithmetic & \texttt{+}, \texttt{-}, \texttt{*}, \texttt{/}, \texttt{\%} & Basic arithmetic \\
Data       & \texttt{var}, \texttt{missing}, \texttt{missing\_some} & Variable access \\
\bottomrule
\end{tabularx}
\end{table}

\subsection{Variable paths}

Variable paths use dot notation to access nested fields in the invoice record.
The root object is the validated invoice conforming to
\texttt{invoice.schema.json}.

\begin{lstlisting}[basicstyle=\ttfamily\small,caption={Example variable paths for invoice data.}]
totals.gross           # Invoice gross amount
totals.vat_rate_percent  # VAT rate as percentage
seller.tax_id          # Seller VAT registration number
seller.country         # Seller country code
buyer.country          # Buyer country code
metadata.issue_date    # Invoice issue date
line_items.0.description  # First line item description
\end{lstlisting}

\subsection{Named helper invocation}

When a predicate requires complex logic beyond JSON-Logic capabilities,
the engine invokes a named helper. The syntax is:

\begin{lstlisting}[basicstyle=\ttfamily\small,caption={Named helper invocation in rule predicate.}]
{
"predicate": {
"type": "named_helper",
"name": "check_vat_registration_validity",
"args": {
"tax_id": {"var": "seller.tax_id"},
"country": {"var": "seller.country"}
}
}
}
\end{lstlisting}

\section{Error handling specification}

The module \textbf{must} handle errors gracefully without crashing the
parent invoice processing pipeline.

\subsection{Error codes}

\begin{table}[htbp]
\centering
\caption{VAT module error codes.}
\label{tab:vat-errors}
\footnotesize
\begin{tabularx}{\linewidth}{@{}l l X@{}}
\toprule
\textbf{Code} & \textbf{Severity} & \textbf{Description} \\ \midrule
\texttt{VAT\_E001} & ERROR & Invalid VAT registration format \\
\texttt{VAT\_E002} & ERROR & Unknown country code \\
\texttt{VAT\_E003} & ERROR & Invalid VAT rate (negative or $>$100\%) \\
\texttt{VAT\_E004} & ERROR & JSON-Logic parse failure \\
\texttt{VAT\_W001} & WARNING & Authority API unavailable (fallback to format check) \\
\texttt{VAT\_W002} & WARNING & Activity code not in lookup table \\
\texttt{VAT\_I001} & INFO & Rate changed since last validation \\
\bottomrule
\end{tabularx}
\end{table}

\subsection{Fallback behaviour}

When an error occurs:

\begin{enumerate}
\item Log the error with full context (invoice ID, field values, stack trace)
\item Set the rule result to \textsc{hold} (never silently pass)
\item Populate \texttt{checklist[].error\_code} with the appropriate code
\item Continue processing remaining rules (fail-open for diagnostics)
\item Return the partial checklist to the caller
\end{enumerate}

\section{Performance requirements}

\begin{table}[htbp]
\centering
\caption{VAT module performance requirements.}
\label{tab:vat-perf}
\footnotesize
\begin{tabularx}{\linewidth}{@{}l r X@{}}
\toprule
\textbf{Metric} & \textbf{Target} & \textbf{Measurement} \\ \midrule
p95 latency per invoice & $<$ 100 ms & Excluding authority API calls \\
p95 latency with API    & $<$ 2000 ms & Including authority validation \\
Memory per invoice      & $<$ 10 MB & Peak allocation during processing \\
Throughput              & $>$ 100/sec & Single-threaded, no API calls \\
\bottomrule
\end{tabularx}
\end{table}

\section{Implementation checklist}

Developers \textbf{must} verify the following before marking the module
complete:

\begin{itemize}
\item[\texttt{[ ]}] All four helper functions implemented with correct signatures
\item[\texttt{[ ]}] Unit tests achieve $\geq$90\% line coverage
\item[\texttt{[ ]}] All country-specific regex patterns validated against test data
\item[\texttt{[ ]}] JSON-Logic parser handles all operators in Table~\ref{tab:jsonlogic-ops}
\item[\texttt{[ ]}] Error codes match Table~\ref{tab:vat-errors}
\item[\texttt{[ ]}] Performance meets targets in Table~\ref{tab:vat-perf}
\item[\texttt{[ ]}] Integration test with real invoice records from each country
\item[\texttt{[ ]}] Security review completed (input validation, injection prevention)
\end{itemize}
\chapter{Testing Specification}\label{ch:testing}

This chapter specifies the testing requirements for the Arabic AP invoice
processing pipeline, including test cases, coverage targets, and CI/CD
integration.

\section{Test coverage requirements}

\begin{table}[htbp]
\centering
\caption{Arabic AP test coverage targets.}
\label{tab:arabic-coverage}
\footnotesize
\begin{tabularx}{\linewidth}{@{}l r r X@{}}
\toprule
\textbf{Component} & \textbf{Line \%} & \textbf{Branch \%} & \textbf{Notes} \\ \midrule
OCR extractor & 85 & 75 & \\
VAT rules engine & 90 & 85 & Critical path \\
Invoice validator & 90 & 85 & Schema + business rules \\
AP workflow engine & 80 & 70 & \\
\textbf{Overall minimum} & \textbf{85} & \textbf{75} & Enforced in CI \\
\bottomrule
\end{tabularx}
\end{table}

\section{Test fixtures}

Fixtures are located in \texttt{tests/fixtures/arabic/}.

\subsection{Invoice fixtures}

\begin{table}[htbp]
\centering
\caption{Required invoice fixture categories.}
\footnotesize
\begin{tabularx}{\linewidth}{@{}l l X@{}}
\toprule
\textbf{Category} & \textbf{Count} & \textbf{Coverage} \\ \midrule
Standard invoices & 20 & 5 per country (SA, EG, BH, QA) \\
Edge cases & 15 & Missing fields, unusual formats \\
Invalid invoices & 10 & Schema violations, invalid TINs \\
Credit notes & 5 & Negative amounts, adjustments \\
Multi-currency & 5 & Non-local currency invoices \\
\bottomrule
\end{tabularx}
\end{table}

\subsection{OCR fixtures}

OCR fixtures cover all confidence thresholds:
\begin{itemize}
\item High quality ($>$85\%): 10 samples
\item Medium quality (55--85\%): 10 samples
\item Low quality (30--55\%): 5 samples
\item Very low quality ($<$30\%): 5 samples
\end{itemize}

\section{Unit test cases}

\subsection{Invoice extraction}

\begin{tabularx}{\linewidth}{@{}l l X@{}}
\toprule
\textbf{ID} & \textbf{Priority} & \textbf{Description} \\ \midrule
AP-EXT-001 & P1 & Extract all fields from standard Saudi invoice \\
AP-EXT-002 & P1 & Extract all fields from standard Egyptian invoice \\
AP-EXT-003 & P1 & Extract all fields from standard Bahrain invoice \\
AP-EXT-004 & P1 & Extract all fields from Qatar payment instruction \\
AP-EXT-005 & P1 & Handle missing optional fields gracefully \\
AP-EXT-006 & P2 & Handle multi-page invoice \\
AP-EXT-007 & P2 & Handle poor quality scan ($<$55\% confidence) \\
AP-EXT-008 & P2 & Handle Arabic-only invoice \\
AP-EXT-009 & P2 & Handle Arabic+English bilingual invoice \\
AP-EXT-010 & P3 & Handle unusual date formats \\
\bottomrule
\end{tabularx}

\subsection{VAT validation}

\begin{tabularx}{\linewidth}{@{}l l X@{}}
\toprule
\textbf{ID} & \textbf{Priority} & \textbf{Description} \\ \midrule
AP-VAT-001 & P1 & Validate 15-digit Saudi TIN format \\
AP-VAT-002 & P1 & Validate 9-digit Egyptian TRN format \\
AP-VAT-003 & P1 & Calculate VAT amount matches invoice \\
AP-VAT-004 & P1 & Apply correct VAT rate by country \\
AP-VAT-005 & P1 & Identify reverse-charge transactions \\
AP-VAT-006 & P2 & Handle zero-rated supplies \\
AP-VAT-007 & P2 & Handle exempt supplies \\
AP-VAT-008 & P2 & Apply withholding tax rules \\
AP-VAT-009 & P3 & Handle credit notes correctly \\
AP-VAT-010 & P3 & Handle inter-GCC transactions \\
\bottomrule
\end{tabularx}

\section{Integration test suite}

\begin{lstlisting}[basicstyle=\ttfamily\small]
pytest tests/integration/arabic/ -v --tb=short
\end{lstlisting}

Expected duration: 15 minutes.

\subsection{End-to-end scenarios}

\begin{enumerate}
\item \textbf{Happy path}: Invoice scan $\rightarrow$ OCR $\rightarrow$ extraction
$\rightarrow$ VAT validation $\rightarrow$ approval workflow
\item \textbf{Manual review}: Low-confidence extraction triggers human review
\item \textbf{VAT hold}: Rule failure triggers hold and escalation
\item \textbf{Multi-country}: Invoices from all four countries processed correctly
\end{enumerate}

\section{Performance benchmarks}

All p95 targets in Table~\ref{tab:ap-perf-targets} are measured on the
\emph{reference hardware profile} in Section~\ref{sec:ap-perf-hardware}.
CI enforces each target by failing the performance regression job when
the measured p95 exceeds the target by more than 20\,\% on three
consecutive runs on identical hardware.

\begin{table}[htbp]
\centering
\caption{Arabic AP performance targets on reference hardware.}
\label{tab:ap-perf-targets}
\footnotesize
\begin{tabularx}{\linewidth}{@{}l r r X@{}}
\toprule
\textbf{Operation} & \textbf{p95 Target} & \textbf{Throughput} & \textbf{Notes} \\ \midrule
OCR extraction & 3,000 ms & 200/hr & Per page \\
Field extraction (LLM) & 5,000 ms & 100/hr & Per invoice \\
VAT rule evaluation & 100 ms & 1,000/sec & No external API \\
Full pipeline E2E & 10,000 ms & 50/hr & Including all steps \\
\bottomrule
\end{tabularx}
\end{table}

\subsection{Reference hardware profile}
\label{sec:ap-perf-hardware}

\begin{description}
\item[GPU.] 1$\times$ NVIDIA A100 40\,GB SXM4 \emph{or} L40S 48\,GB.
Shared with Simula and TB; reservation is not required for
single-invoice latency targets.
\item[OCR engine.] Tesseract 5.3.x with Arabic language pack
(\texttt{ara}, \texttt{ara\_old}), CPU-only. A 2-page invoice
at 300\,dpi is the reference document size.
\item[LLM.] \texttt{Qwen2.5-14B-Instruct} served via vLLM 0.6.x,
\texttt{max\_model\_len=8192}, 14-field extraction prompt
($\sim$1,200 input tokens, $\sim$400 output tokens).
\item[CPU / RAM.] 8~vCPU, 32\,GB RAM host.
\item[Storage.] NVMe SSD; invoice PDFs streamed from object storage
with 2\,Gbps link.
\item[Network.] Co-located vLLM; RTT $<$0.5\,ms. For remote LLM
endpoints, add measured RTT to the field-extraction p95.
\end{description}

\noindent
Running on consumer-grade GPUs (e.g.\ T4 15\,GB) is not supported for
the 14B model; a fallback profile using
\texttt{Qwen2.5-7B-Instruct} is documented in
\path{docs/arabic/runbooks/perf-fallback.md} with correspondingly
relaxed p95 targets (field extraction $\le$8,000\,ms).

\section{Security tests}

\subsection{Prompt injection}

\begin{tabularx}{\linewidth}{@{}l X@{}}
\toprule
\textbf{Test} & \textbf{Description} \\ \midrule
AP-SEC-001 & Malicious text in invoice description field \\
AP-SEC-002 & JSON injection via seller name \\
AP-SEC-003 & Unicode control characters in OCR output \\
AP-SEC-004 & System prompt extraction attempt \\
\bottomrule
\end{tabularx}

\section{CI/CD integration}

\begin{lstlisting}[basicstyle=\ttfamily\small,caption={GitHub Actions workflow excerpt.}]
arabic-tests:
runs-on: ubuntu-latest
steps:
- uses: actions/checkout@v4
- name: Run Arabic AP tests
run: |
pytest tests/unit/arabic/ -v --cov=src/arabic
pytest tests/integration/arabic/ -v
- name: Check coverage
run: |
coverage report --fail-under=85
\end{lstlisting}

\section{Acceptance criteria}

Tests pass when:
\begin{enumerate}
\item All P1 test cases pass (100\%)
\item P2 test cases pass ($\geq$95\%)
\item P3 test cases pass ($\geq$90\%)
\item Line coverage $\geq$85\%
\item Performance within targets
\item No critical security findings
\end{enumerate}
\chapter{OCR and Translation Calibration Pilot Design}\label{ch:ocr-pilot}

This chapter specifies the pilot study design for calibrating OCR confidence
thresholds \emph{and} the Arabic$\to$English translation quality gate,
addressing the review findings that both the OCR thresholds (30\,\%/55\,\%)
and the Levenshtein distance gate ($\le$0.30) are calibrated on three
documents. Both calibrations are run against the same pilot corpus to avoid
duplicate ground-truth effort.

\section{Objective}

The pilot aims to:
\begin{enumerate}
\item Determine empirically-validated OCR confidence thresholds per country
\item Establish confidence intervals for threshold boundaries
\item Identify document types that require special handling
\item Create a baseline for ongoing threshold monitoring
\item Calibrate the Arabic$\to$English translation quality gate
(Levenshtein distance threshold) against the same pilot corpus
\end{enumerate}

\section{Pilot design}

\subsection{Sample size}

\begin{table}[htbp]
\centering
\caption{Pilot sample size by country.}
\label{tab:pilot-sample}
\footnotesize
\begin{tabularx}{\linewidth}{@{}l r r X@{}}
\toprule
\textbf{Country} & \textbf{Minimum} & \textbf{Target} & \textbf{Stratification} \\ \midrule
Saudi Arabia (SA) & 50 & 75 & By invoice type (tax, commercial) \\
Egypt (EG) & 50 & 75 & By scan quality (photo, scan, PDF) \\
Bahrain (BH) & 50 & 75 & By document language (AR, EN, mixed) \\
Qatar (QA) & 50 & 75 & By amount range (small, medium, large) \\
\textbf{Total} & \textbf{200} & \textbf{300} & \\
\bottomrule
\end{tabularx}
\end{table}

The sample size of 50 per country provides:
\begin{itemize}
\item 95\% confidence interval width of $\pm$7\% for threshold estimates
\item Statistical power of 80\% to detect 15\% differences between countries
\item Sufficient stratification for 3--4 document types per country
\end{itemize}

\subsection{Document selection}

Documents are selected by:
\begin{enumerate}
\item Random sampling from the last 6 months of invoices
\item Stratified by document type (tax invoice, commercial invoice, payment
instruction)
\item Stratified by scan quality (high: $>$300 DPI, medium: 150--300 DPI,
low: $<$150 DPI)
\item Excluding documents with known data quality issues
\end{enumerate}

\subsection{Ground truth creation}

For each pilot document:
\begin{enumerate}
\item Two independent human reviewers manually transcribe all fields
in the source language (Arabic) and the required English fields
\item Inter-rater agreement calculated (target: Cohen's $\kappa > 0.9$)
for categorical fields; character error rate (CER) $\le$1\,\%
on free-text fields
\item Disagreements are resolved by a third reviewer whose decision is
final. The resolution process is:
\begin{enumerate}
\item Third reviewer inspects both transcriptions side-by-side
with the source document.
\item Third reviewer selects one of the two transcriptions
\emph{or} provides a new transcription; the choice and
rationale are logged to
\path{pilot/disputes/<doc-id>.json}.
\item The third reviewer's verdict is the ground truth; the
pilot lead reviews the top 10\,\% of disputes monthly
to detect systematic bias.
\end{enumerate}
\item Final ground truth stored in structured JSON
(\path{pilot/ground-truth/<country>/<doc-id>.json}), validated
against \texttt{docs/schema/arabic/ocr-result.schema.json}
\end{enumerate}

\section{Threshold calibration methodology}

\subsection{OCR evaluation}

Each document is processed by the OCR pipeline:
\begin{enumerate}
\item Record per-field OCR confidence scores
\item Record per-document aggregate confidence
\item Compare OCR output to ground truth
\item Calculate field-level accuracy (character error rate)
\end{enumerate}

\subsection{Threshold determination}

For each country, determine optimal thresholds using:

\begin{lstlisting}[basicstyle=\ttfamily\small,caption={Threshold calibration algorithm.}]
def calibrate_thresholds(pilot_data, target_precision=0.95):
"""
Calibrate OCR confidence thresholds.

Parameters:
pilot_data: List of (confidence, accuracy) pairs
target_precision: Minimum acceptable precision (default 95%)

Returns:
high_threshold: Above this, auto-accept
low_threshold: Below this, reject
"""
# Sort by confidence descending
sorted_data = sorted(pilot_data, key=lambda x: -x.confidence)

# Find high_threshold: highest confidence where
# precision >= target for all docs above threshold
high_threshold = find_precision_cutoff(
sorted_data, target_precision
)

# Find low_threshold: confidence below which
# accuracy < 50% (not worth manual review)
low_threshold = find_accuracy_floor(sorted_data, floor=0.50)

return high_threshold, low_threshold
\end{lstlisting}

\paragraph{Justification of \texttt{target\_precision = 0.95}.}
The AP workflow routes any invoice below the high threshold to
two-person manual review (Chapter~5), which is tolerant of 5\,\%
false-positive auto-accepts from a \emph{business-risk} standpoint:
the downstream VAT engine (Chapter~4) independently validates
tax-ID and amount fields, and the reconciliation step catches
posting errors before payment is released. The 97\,\% soundness
threshold reported in the referenced OCR evaluation paper is
appropriate for \emph{unsupervised} pipelines; because our pipeline
has a manual review backstop, 95\,\% provides the right
precision/recall trade-off (empirically validated in the pilot
success criteria below). A sensitivity analysis at 93\,\%, 95\,\%,
and 97\,\% must be produced as part of the pilot report; the final
threshold is the largest value in \{0.93, 0.95, 0.97\} for which
the success criteria in Section~\ref{sec:pilot-success} hold.

\subsection{Translation quality gate calibration}
\label{sec:pilot-translation}

Using the same pilot corpus, the Arabic$\to$English translation
quality gate (currently Levenshtein distance $\le$0.30 against the
model's re-translation, calibrated on three documents) is
re-calibrated:

\begin{enumerate}
\item For each ground-truthed pilot document, run the
Arabic$\to$English$\to$Arabic round-trip translation and
compute the normalized Levenshtein distance against the
source Arabic.
\item Compute the field-level translation accuracy against the
human ground truth (exact match on numeric fields, BLEU-4
$\ge$0.6 on free-text fields).
\item Determine the Levenshtein threshold $\tau^{*}$ that
maximises the F1 between ``round-trip gate passed'' and
``field-level translation accurate'' across the 200 pilot
documents.
\item Report the per-country threshold and the 95\,\% bootstrap
CI ($\tau^{*}_{\mathrm{lower}}, \tau^{*}_{\mathrm{upper}}$).
Deploy using the conservative upper bound.
\end{enumerate}

\noindent
The translation calibration adds no additional ground-truth effort
(the human transcriptions already contain the English fields) and
extends the pilot timeline by one week (added to
Section~\ref{sec:pilot-timeline}).

\subsection{Confidence interval calculation}

Bootstrap sampling (1,000 iterations) determines confidence intervals:
\begin{itemize}
\item Report 95\% CI for high threshold
\item Report 95\% CI for low threshold
\item Use upper bound of CI for conservative deployment
\end{itemize}

\section{Success criteria}\label{sec:pilot-success}

The pilot succeeds when all of the following are true:
\begin{enumerate}
\item $\geq$200 documents processed (50 per country minimum)
\item Inter-rater agreement $\kappa > 0.9$ (categorical fields)
and CER $\le$1\,\% (free-text fields)
\item 95\,\% CI width $\leq \pm 5\,\%$ for all thresholds
\item High threshold achieves $\geq$95\,\% precision
\item Low threshold correctly rejects $\geq$90\,\% of unusable documents
\item Translation gate F1 $\geq$0.85 at the selected Levenshtein
threshold (see Section~\ref{sec:pilot-translation})
\end{enumerate}

\section{Fallback procedure}

If the pilot fails success criteria:
\begin{enumerate}
\item Increase sample size by 50\% and repeat
\item If still failing, segment by document type and calibrate separately
\item If per-type calibration fails, default to manual review for all
documents from that country until more data is collected
\end{enumerate}

\section{Timeline}\label{sec:pilot-timeline}

\begin{table}[htbp]
\centering
\caption{OCR and translation pilot timeline.}
\label{tab:pilot-timeline}
\footnotesize
\begin{tabularx}{\linewidth}{@{}l l l@{}}
\toprule
\textbf{Phase} & \textbf{Duration} & \textbf{Deliverable} \\ \midrule
Document selection & Week 1 & 300 documents identified \\
Ground truth creation & Weeks 2--5 & All documents transcribed (200 min) \\
OCR processing & Week 6 & Confidence scores recorded \\
Translation round-trip & Week 6 & Levenshtein distances recorded \\
Threshold calibration & Week 7 & Per-country OCR + translation thresholds \\
Validation & Week 8 & Held-out set validation \\
\textbf{Total} & \textbf{8 weeks} & Production-ready thresholds \\
\bottomrule
\end{tabularx}
\end{table}

\paragraph{Ground-truth effort justification.}
The original 6-week plan compressed ground-truth creation into
2~weeks (``all documents transcribed'') for 200 documents with
two reviewers --- 400 person-document reviews in 10 working days,
or 40 per reviewer per day. Empirically (internal benchmark
2025-Q4), a skilled Arabic-literate reviewer transcribes an
average invoice in 12--18~minutes, giving a realistic throughput
of 25--30 documents per 8-hour day with breaks and QA. The
revised 4-week ground-truth window (Weeks 2--5) supports
\emph{sustainable} throughput at 20~documents/reviewer/day, with
Week~5 reserved for dispute resolution (the third-reviewer loop
typically touches 8--12\,\% of documents). Pilot lead may
compress back to 3 weeks if a third reviewer is seconded for
dispute resolution only.

\section{Ownership}

\begin{table}[htbp]
\centering
\caption{Pilot roles and responsibilities.}
\label{tab:pilot-roles}
\footnotesize
\begin{tabularx}{\linewidth}{@{}l X@{}}
\toprule
\textbf{Role} & \textbf{Responsibility} \\ \midrule
Pilot Lead & Overall coordination, success criteria validation \\
Data Analyst & Sample selection, threshold calibration \\
QA Reviewers (2) & Ground truth creation, inter-rater agreement \\
ML Engineer & OCR pipeline configuration, metrics collection \\
Product Owner & Go/no-go decision, fallback approval \\
\bottomrule
\end{tabularx}
\end{table}

\section{Post-pilot monitoring}

After threshold deployment:
\begin{enumerate}
\item Monthly accuracy review on random 50-document sample
\item Alert if accuracy drops $>$5\% from pilot baseline
\item Annual recalibration with new 200-document sample
\item Per-country threshold adjustment if document mix changes
\end{enumerate}
\chapter{Project Roadmap}\label{ch:ap-roadmap}

\section*{Scope of this chapter}
\addcontentsline{toc}{section}{Scope of this chapter}

This chapter is the roadmap for the \textbf{Arabic invoice processing}
project. It is a project-level artefact, not a programme plan. Other
SAP OSS Finance projects (Trial Balance, Simula, Regulations) run on
the same platform, workspaces, and backends, but each owns its own
roadmap, business case, and go/no-go decision. This chapter adopts the
MoSCoW prioritisation and sequence-based critical path that the four
projects share so that artefacts can be compared at a glance; it does
not assert coupling between their go-live moments.

\paragraph{Sequence, not dates.}
The milestones and dependency graph below express \emph{order} and
\emph{prerequisite} relationships, not calendar dates. This is
deliberate. Implementation teams must know what depends on what, and
where the schedule can absorb slippage (the Should/Could items); a
date column would harden targets into commitments they were never
meant to be. Dates, when set, live in the project steering packet and
the ticketing system, not in this specification.

\section{MoSCoW classification}\label{sec:ap-moscow}

Every milestone below carries one of four priorities. The mapping is
contractual: a milestone's priority determines what happens when it is
dropped, de-scoped, or slips.

\begin{table}[htbp]
\centering
\caption{MoSCoW prioritisation (shared across the four SAP OSS Finance
projects).}
\label{tab:ap-moscow}
\footnotesize
\begin{tabularx}{\linewidth}{@{}l l X@{}}
\toprule
\textbf{Priority} & \textbf{Meaning} & \textbf{Allowance} \\ \midrule
\textbf{Must} (\textbf{M}) & Non-negotiable for v1.0. Without this milestone, the release does not ship. & None. Missing a Must halts the release; see the kill-switch in \cref{sec:ap-kill-switch}. \\
\textbf{Should} (\textbf{S}) & Important but not vital. Ship without it only with sponsor sign-off; document the deviation in the business-case refresh. & May slip by one rollout wave; two consecutive slips escalate to GFAIC. \\
\textbf{Could} (\textbf{C}) & Desirable enhancement. Included if Must and Should items are on track. & May be dropped to v1.1 at project steering's discretion. \\
\textbf{Won't} (\textbf{W}) & Explicitly out of v1.0 scope. Listed so the boundary is legible. & Not eligible for v1.0 planning at all. \\
\bottomrule
\end{tabularx}
\end{table}

\noindent
The critical path of this project is the chain of Must items in
\cref{tab:ap-milestones}; the Should and Could items are where the
schedule and scope can absorb delivery risk without a release slip.

\section{Dependency graph}\label{sec:ap-dependency-graph}

\begin{figure}[htbp]
\centering
\begin{tikzpicture}[
node distance=5mm and 10mm,
must/.style={draw,rounded corners=2pt,align=center,font=\scriptsize,
fill=apred!20,draw=apred!70,
minimum height=7mm,minimum width=32mm},
should/.style={draw,rounded corners=2pt,align=center,font=\scriptsize,
fill=apblue!12,draw=apblue!50,
minimum height=7mm,minimum width=32mm},
could/.style={draw,rounded corners=2pt,align=center,font=\scriptsize,
fill=apgreen!12,draw=apgreen!50,dotted,
minimum height=7mm,minimum width=32mm},
dep/.style={-{Latex[length=2mm]},thick},
soft/.style={-{Latex[length=2mm]},thick,dashed}
]
\node[must] (m01) {M01 AP-ENG-001 \\ \scriptsize OCR pipeline};
\node[must,below=of m01] (m02) {M02 AP-ENG-002 \\ \scriptsize VAT engine};
\node[must,below=of m02] (m03) {M03 AP-PILOT-PREP \\ \scriptsize 200 pilot docs};
\node[must,right=14mm of m02] (m04) {M04 REG-INTG-AP-001 \\ \scriptsize identity wiring};
\node[must,right=14mm of m04] (m05) {M05 AP-PILOT-001 \\ \scriptsize thresholds set};
\node[must,below=of m05] (m06) {M06 AP-BETA-001 \\ \scriptsize Saudi beta};
\node[must,below=of m06] (m07) {M07 AP-PARALLEL-001 \\ \scriptsize parallel-run gate};
\node[must,below=of m07] (m08) {M08 AP-GA-001 \\ \scriptsize Saudi GA};
\node[should,right=14mm of m08] (s01) {S01 AP-WAVE-002 \\ \scriptsize Egypt};
\node[should,above=of s01] (s02) {S02 AP-WAVE-003 \\ \scriptsize Bahrain};
\node[should,above=of s02] (s03) {S03 AP-WAVE-004 \\ \scriptsize Qatar};
\node[could,above=of s03] (c01) {C01 AP-V2-EXPLORE \\ \scriptsize multilingual};

\draw[dep] (m01) -- (m02);
\draw[dep] (m01) -- (m03);
\draw[dep] (m02) -- (m04);
\draw[dep] (m03) -- (m05);
\draw[dep] (m04) -- (m05);
\draw[dep] (m05) -- (m06);
\draw[dep] (m06) -- (m07);
\draw[dep] (m07) -- (m08);
\draw[dep] (m08) -- (s01);
\draw[dep] (s01) -- (s02);
\draw[dep] (s02) -- (s03);
\draw[soft] (s03) -- (c01);
\end{tikzpicture}
\caption{Arabic AP dependency graph. Red nodes are Must (critical
path); blue are Should (country-wave allowance); dotted green are
Could (discretionary). Solid arrows are hard prerequisites; dashed
arrows indicate conditional starts (the target is unlocked when the
upstream completes \emph{and} the multilingual volume trigger in
\cref{sec:ap-v2-trigger} fires).}
\label{fig:ap-dependency-graph}
\end{figure}

\section{Project milestones}\label{sec:ap-milestones}

\begin{table}[htbp]
\centering
\caption{Arabic AP v1.0 project milestones in dependency order.
Priority: M=Must (critical path), S=Should, C=Could, W=Won't.}
\label{tab:ap-milestones}
\footnotesize
\begin{tabularx}{\linewidth}{@{}l l l l X@{}}
\toprule
\textbf{ID} & \textbf{Prereq} & \textbf{Pri.} & \textbf{Ticket} & \textbf{Milestone} \\ \midrule
M01 & ---       & M & AP-ENG-001         & OCR pipeline (Tesseract + Arabic model) passes the fixtures in \cref{ch:testing}. \\
M02 & M01       & M & AP-ENG-002         & VAT engine unit tests pass (\cref{ch:vat}); country-configurable rates.\\
M03 & M01       & M & AP-PILOT-PREP      & 200 pilot documents sampled and ground-truth annotation complete (\cref{ch:ocr-pilot}). \\
M04 & M02       & M & REG-INTG-AP-001    & Identity-attribution wiring (reserve/complete) live in UAT. \\
M05 & M03, M04  & M & AP-PILOT-001       & Pilot complete; OCR and translation thresholds calibrated per country. \\
M06 & M05       & M & AP-BETA-001        & Beta cycle: Saudi Arabia only, two-person manual review on top of AI. \\
M07 & M06       & M & AP-PARALLEL-001    & Beta parallel-run gate passes. \\
M08 & M07       & M & AP-GA-001          & GA sign-off; Saudi Arabia goes live. \\
M09 & M08       & M & AP-BC-REFRESH      & Business-case refresh and compliance attestation (first cycle after GA; then annual). \\
S01 & M08       & S & AP-WAVE-002        & Rollout wave 2: Egypt onboarded via country-params. \\
S02 & S01       & S & AP-WAVE-003        & Rollout wave 3: Bahrain onboarded. \\
S03 & S02       & S & AP-WAVE-004        & Rollout wave 4: Qatar onboarded (full v1.0 scope). \\
C01 & S03 + \cref{sec:ap-v2-trigger} trigger & C & AP-V2-EXPLORE & v2.0 multilingual expansion scoped (Hebrew, Farsi) if volume justifies. \\
C02 & M08       & C & AP-PROMPT-TUNE     & Quarterly translation-prompt retraining from accept/edit telemetry. \\
W01 & ---       & W & AP-V1-MULTILINGUAL & Non-Arabic-script languages are \emph{not} in v1.0 scope. \\
W02 & ---       & W & AP-V1-ONE-PERSON   & One-person review is \emph{not} in v1.0 scope; relaxed only after full-rollout gate sustained. \\
\bottomrule
\end{tabularx}
\end{table}

\noindent
The Must chain (M01 $\to$ M02/M03 $\to$ M04 $\to$ M05 $\to$ M06 $\to$
M07 $\to$ M08 $\to$ M09) is the critical path. The Should items
(S01--S03) are the primary \emph{schedule allowance}: if Must items
run hot, country waves can slip without affecting v1.0 sign-off. The
Could items are the \emph{scope allowance}: they may be cut to v1.1
without renegotiation.

\section{Launch tiers and exit gates}\label{sec:ap-launch-tiers}

The project walks the four-tier ladder shared across SAP OSS Finance
projects. A tier is exited only when the named exit gate is satisfied;
slippage requires sponsor sign-off.

\begin{table}[htbp]
\centering
\caption{Arabic AP launch tiers. The exit gate defines the MoSCoW
items that must be closed to exit the tier.}
\label{tab:ap-launch-tiers}
\footnotesize
\begin{tabularx}{\linewidth}{@{}l l X@{}}
\toprule
\textbf{Tier} & \textbf{Scope} & \textbf{Exit gate (references MoSCoW IDs)} \\ \midrule
Alpha & Engineering only, DEV + UAT, fixture documents only. & M01, M02, M03 closed; all acceptance criteria in \cref{ch:testing} pass on the pilot fixture; reproducibility of OCR and VAT outputs confirmed on two consecutive runs. \\
Beta  & One country (Saudi Arabia), two-person manual review on top of AI. & M04, M05, M06, M07 closed; pilot success criteria (\cref{sec:pilot-success}) met; parallel-run gate passes for one full cycle; no open SEV1/SEV2; audit coverage $=$ 100\,\% of in-scope actions. \\
GA (limited) & Three countries (Saudi Arabia plus the next two waves). & M08 closed; Beta gate sustained for two consecutive cycles; S01 closed; country-onboarding runbook exercised end-to-end; kill-switch green; per-country thresholds calibrated. \\
Full rollout & All four countries in v1.0 scope (SA, EG, BH, QA). & S02 and S03 closed; two quarterly reviews sustained; M09 complete; business-case refresh complete; two-person review may be relaxed to one-person for documents above the high threshold. \\
\bottomrule
\end{tabularx}
\end{table}

\section{Kill-switch conditions}\label{sec:ap-kill-switch}

The kill-switch is a pre-committed rulebook: a tripped condition enters
the action column without further debate. GFAIC and sponsor are
notified within 24 hours. The conditions are evaluated continuously
(on every cycle) and do not carry dates.

\begin{table}[htbp]
\centering
\caption{Arabic AP kill-switch.}
\label{tab:ap-kill-switch}
\footnotesize
\begin{tabularx}{\linewidth}{@{}l l X@{}}
\toprule
\textbf{Condition} & \textbf{Evaluation} & \textbf{Action} \\ \midrule
Pilot F1 on the translation gate $<$\,0.70 on three of four countries & on M05 close and any recalibration & \textbf{Redesign}: narrow v1.0 to best-performing countries only; full rollout paused. \\
Posted-VAT error rate $>$\,0.5\,\% over any cycle in Beta or GA & per cycle & \textbf{Halt onboarding}; incident-review before next wave; re-calibrate thresholds. \\
Two-person review sustains $>$\,20\,\% edit rate for three consecutive cycles in GA & per cycle & \textbf{Redesign prompt layer}; commentary drafts gated behind flag until edit rate $\le$\,15\,\%. \\
Audit-coverage gap (missed \texttt{reserve/complete} pair) detected in production & continuous & \textbf{Pause GA onboarding}; remediate before next wave. \\
Any MAS MGF or country-VAT regulator finding attributable to AP specifically & continuous & \textbf{Halt rollout}; remediate before next wave. \\
Shared-platform dependency in \cref{tab:ap-shared-deps} declared red by its owner project & continuous & \textbf{Pause downstream Must}; re-plan at next steering. \\
\bottomrule
\end{tabularx}
\end{table}

\section{Shared-platform dependencies}\label{sec:ap-shared-deps}

This project depends on four shared-platform components. Each is owned
by a sibling project (Regulations) or by shared infra, and is tracked
here only as a dependency, not a commitment this project can
unilaterally adjust. Sequence is expressed in terms of the Must item
that first consumes the dependency, not in calendar time.

\begin{table}[htbp]
\centering
\caption{Shared-platform dependencies.}
\label{tab:ap-shared-deps}
\footnotesize
\begin{tabularx}{\linewidth}{@{}l l l X@{}}
\toprule
\textbf{Component} & \textbf{Owner} & \textbf{First consumed by} & \textbf{Impact if late} \\ \midrule
Audit service (reserve/complete, immutable storage) & Regulations project & M04 (REG-INTG-AP-001) & No Beta gate possible without this; AP GA slips in lockstep with the audit service. \\
Unified schema registry (\path{docs/schema/}) & Shared infra & M01 & Already \texttt{migration.status=complete}; no remaining risk. \\
vLLM hosting (\texttt{vllm-only} policy) & Shared platform & M02 & Graceful-degradation ladder defined in the Regulations spec provides fallback; slip tolerable. \\
Tesseract + Arabic OCR model & Shared platform & M01 & No substitute available; slip blocks all subsequent Must items. \\
\bottomrule
\end{tabularx}
\end{table}

\noindent
Shared-platform dependencies are surfaced at monthly steering. If any
upstream owner declares \textbf{amber} or \textbf{red} on a component
in \cref{tab:ap-shared-deps}, the affected downstream Must milestone
is re-planned at the following steering.

\section{v2.0 expansion trigger}\label{sec:ap-v2-trigger}

The Could item C01 (AP-V2-EXPLORE) is unlocked only when all of the
following are observed on production data from S03 (Qatar onboarded):

\begin{itemize}
\item aggregate non-Arabic-script document volume $\ge\,15$\,\% of
total Arabic AP volume for two consecutive quarters;
\item translation-gate F1 on Arabic above the Beta threshold
sustained for the same window;
\item business-case refresh (M09) explicitly budgets the v2.0
programme.
\end{itemize}

\noindent
If any of the three conditions fails, C01 remains in the Could column
and is re-evaluated at the next business-case refresh.

\section{Governance and business-case cadence}\label{sec:ap-governance}

\paragraph{Project governance.}
The Arabic AP project reports monthly to a project steering (AP
Product Owner, Tax Controller, AI Engineer) and quarterly to the
cross-project GFAIC. Monthly steering accepts or rejects amber/red
signals on Must items; GFAIC reviews compliance posture and any
cross-project dependencies.

\paragraph{Business-case refresh (M09).}
The business case that authorises this project is refreshed on the
following rhythm:

\begin{itemize}
\item \textbf{Mandatory refresh}: at the first cycle after GA, and
annually thereafter (ticket AP-BC-REFRESH, priority M in
\cref{tab:ap-milestones}). Refresh packet includes country
volume assumptions, unit economics,
FTE leverage realised vs.\ forecast, and a deviation statement
against the baseline.
\item \textbf{Trigger-based refresh}: any of
\begin{itemize}
\item country monthly volume changes by $\ge\,20$\,\% from
forecast (either direction);
\item per-invoice run cost changes by $\ge\,30$\,\% (OCR
throughput, vLLM licensing, GPU price);
\item new country-specific VAT regulation or MAS MGF pillar
materially affecting scope;
\item any kill-switch in \cref{tab:ap-kill-switch} is tripped.
\end{itemize}
\item \textbf{Owner}: AP Product Owner. Approval routed through the
project steering and filed at
\path{docs/arabic/business-case/refresh-history.yaml}.
\end{itemize}
\chapter{Agent Implementation Instructions}\label{ch:agent-impl}

\section{Purpose and audience}
This chapter is the self-contained, intent-based instruction set for
implementing the \emph{Arabic AP invoice processing} domain inside the
unified SAP OSS platform. It is written for LLM-driven and human
implementation agents who will extend existing services to realise the
specification. It assumes only this PDF; every platform decision,
integration contract, and governance obligation the Arabic workstream
depends on is restated here so that no companion document is required.
Existing chapters~\ref{ch:overview}--\ref{ch:ap-roadmap} remain the
normative source for the Arabic domain itself; this chapter tells
future agents \emph{how} to build what those chapters describe,
\emph{where} to put it in the codebase, and \emph{how} it plugs into
the wider platform.

% !TEX root = ../master.tex
% =============================================================================
% MASTER PLATFORM CORE: Shared Intent-Based Instructions (Union)
% World-Class Version: HITL Gates, Process Integrity, and Sequence Logic
% =============================================================================

\section{Rules of Engagement for AI Agents}
\label{sec:agent-rules}

To ensure safety and architectural integrity, all implementation agents must adhere to the following Prime Directives:

\begin{enumerate}[label=\textsc{rule}-\arabic*:, leftmargin=2cm]
\item \textbf{Credential Protection.} Never log, print, or commit secrets. Use the \texttt{.secrets.example} pattern for all new integrations.
\item \textbf{Context Efficiency.} Prioritize surgical edits. Request only the minimum required file ranges to satisfy the task intent.
\item \textbf{Test-Driven Implementation.} No code is complete without a corresponding fixture in \texttt{tests/fixtures/} and a passing validation run.
\item \textbf{Self-Correction.} If a schema validation fails, the agent must diagnose the root cause (hallucination vs. schema drift) before re-attempting implementation.
\end{enumerate}

\section{Platform Overview --- One Codebase, One Instance}
\label{sec:agent-platform-overview}

\subsection{Integration Stance}
The SAP OSS platform is delivered as a single codebase that runs as a single logical instance. All cross-cutting concerns are hosted by existing services.

\begin{adrbox}{{One Codebase, One Instance}}
To minimize integration latency and eliminate distributed-system failure modes (e.g., partial partitions), we mandate a single logical instance. All cross-domain orchestration happens via the MCP Gateway layer.
\end{adrbox}

\subsection{Platform Integration Map}
\label{sec:platform-map}

\begin{figure}[htbp]
\centering
\begin{tikzpicture}[
node distance=2.5cm,
auto,
block/.style={rectangle, draw, rounded corners, fill=sapblue!10, text width=6em, text centered, minimum height=3em, font=\scriptsize},
regs/.style={rectangle, draw, rounded corners, fill=sapgrey!20, text width=7em, text centered, minimum height=4em, thick, dotted, font=\scriptsize}
]
\node (arabic) [block] {Arabic AP};
\node (tb) [block, below of=arabic] {Trial Balance};
\node (simula) [block, right of=arabic, node distance=5.5cm] {Simula Training};
\node (regs) [regs, below of=simula] {Regulations (Runtime Wrapper)};

\draw [->, thick] (arabic) -- node [align=left, font=\tiny] {Posted\\Invoices} (tb);
\draw [->, thick, dashed] (simula) -- node [align=left, font=\tiny, near start] {Training\\Examples} (arabic);
\draw [->, thick, dashed] (simula) -- node [align=left, font=\tiny, near start] {Training\\Examples} (tb);
\draw [->, thick] (regs) -| node [font=\tiny, near end, xshift=-1cm] {Governance Gate} (arabic);
\draw [->, thick] (regs) -| node [font=\tiny, near end, xshift=-1cm] {Audit \& Identity} (tb);
\end{tikzpicture}
\caption{Cross-domain interaction and unified data flow map. Detailed governance logic is defined in \cref{reg:ch:mgf}.}
\end{figure}

\section{Transactional Sequence Diagrams}
\label{sec:agent-sequences}

Implementation agents must adhere to the temporal logic defined below. Every cross-domain request MUST be brokered by the Regulations Wrapper.

\subsection{Universal Governance Handshake}
This sequence applies to \textbf{all} platform service calls.

\begin{figure}[htbp]
\centering
\begin{tikzpicture}[x=1cm, y=-0.8cm, font=\scriptsize]
% Lifelines
\draw[thick] (0,0) -- (0,7) node[below] {Persona (UI)};
\draw[thick] (3,0) -- (3,7) node[below] {MCP Gateway};
\draw[thick] (6,0) -- (6,7) node[below] {Regs Wrapper};
\draw[thick] (9,0) -- (9,7) node[below] {Domain Service};

% Interaction
\draw[->, >=latex] (0,1) -- node[above] {1. Request + Identity} (3,1);
\draw[->, >=latex] (3,2) -- node[above] {2. Validate (Pillar 1)} (6,2);
\draw[<-, >=latex] (3,3) -- node[above] {3. Auth/Allow Decision} (6,3);
\draw[->, >=latex] (3,4) -- node[above] {4. Execute Tool} (9,4);
\draw[->, >=latex] (9,5) -- node[above] {5. Log Evidence (Pillar 2)} (6,5);
\draw[<-, >=latex] (0,6) -- node[above] {6. Response + Audit ID} (3,6);

\node[draw, fill=yellow!10, text width=2.5cm, align=center] at (4.5,0.2) {MAS MGF Gate};
\end{tikzpicture}
\caption{Standard platform handshake for governed tool execution. Audit schemas are detailed in \cref{reg:ch:identity}.}
\end{figure}

\subsection{Arabic AP Golden Path}
Temporal flow for invoice processing with Arabic context.

\begin{figure}[htbp]
\centering
\begin{tikzpicture}[x=1cm, y=-0.6cm, font=\tiny]
\draw[thick] (0,0) -- (0,8) node[below] {Intake};
\draw[thick] (2.5,0) -- (2.5,8) node[below] {OCR};
\draw[thick] (5,0) -- (5,8) node[below] {VAT Engine};
\draw[thick] (7.5,0) -- (7.5,8) node[below] {Workflow};
\draw[thick] (10,0) -- (10,8) node[below] {Human (HITL)};

\draw[->] (0,1) -- node[above] {Extract} (2.5,1);
\draw[->] (2.5,2) -- node[above] {Translate} (5,2);
\draw[->] (5,3) -- node[above] {Evaluate} (5,4); % Internal loop
\draw[->] (5,5) -- node[above] {\textsc{Hold} Notification} (10,5);
\draw[<-] (5,6) -- node[above] {Manual Release} (10,6);
\draw[->] (5,7) -- node[above] {Final Posting} (7.5,7);
\end{tikzpicture}
\caption{Arabic AP sequence: Scanned PDF to Release. VAT rules are specified in \cref{ar:ch:vat}.}
\end{figure}

\section{Compute and Resource Alignment}
\label{sec:hardware-profiles}

AI algorithms are calibrated for the reference infrastructure established in the domain chapters:

\begin{table}[htbp]
\centering
\footnotesize
\caption{Infrastructure Reference Profiles.}
\label{tab:hardware-profiles}
\begin{tabularx}{\textwidth}{@{}l X X@{}}
\toprule
\textbf{Workload} & \textbf{Reference Hardware} & \textbf{Performance Target} \\
\midrule
LLM Inference & NVIDIA L40S (48GB) & AWQ 4-bit / p95 < 2s \\
Embeddings & NVIDIA A100 (40/80GB) & FP16 / p99 < 500ms \\
Simula Training & NVIDIA A100 SXM4 cluster & BF16 throughput parity \\
\bottomrule
\end{tabularx}
\end{table}

\section{Human-in-the-Loop (HITL) Critical Gates}
\label{sec:hitl-gates}

The following "Four-Eyes" gates are mandatory. AI agents must implement the "Awaiting Sign-off" state for these actions:
\begin{itemize}
\item \textbf{VAT Release:} No invoice marked as \textsc{hold} can be released without human review (see \cref{ar:ch:ap}).
\item \textbf{Anomaly Override:} Statistical outliers ($3\sigma$) cannot be dismissed by AI; only "Confirmed" or "Triaged" by human (see \cref{tb:ch:controls}).
\item \textbf{Journal Posting:} Agents may \emph{draft} journal entries, but only the Financial Controller persona can \emph{commit} to S/4HANA (see \cref{tb:ch:controls-engine}).
\end{itemize}

\section{Registry Integrity \& Live Status}
\label{sec:registry-integrity}

Before initiating any schema-dependent task, implementation agents must verify the live status of the Unified Schema Registry (see \cref{sim:chap:schema}).
\begin{adrbox}{Registry First}
If \texttt{docs/schema/registry.json} is out of sync with the codebase, the agent must halt and emit a \textsc{schema-mismatch} event.
\end{adrbox}

\section{Master Implementation Manifest (Union)}
Implementation runs in three waves across the entire platform.

\subsection*{Wave 1 --- Foundations}
\begin{itemize}
\item \textbf{Simula:} Implement HANA schema extraction pipeline (MCP discovery). See \cref{sim:chap:extraction}.
\item \textbf{Regulations:} Implement MAS MGF Pillar 1 logic (Tool allow-list). See \cref{reg:ch:mgf}.
\item \textbf{Arabic:} Implement Invoice Intake \& OCR Gateway. See \cref{ar:ch:ai}.
\item \textbf{Trial Balance:} Implement identity-attribution wiring. See \cref{reg:ch:identity}.
\end{itemize}

\subsection*{Wave 2 --- Core Domain Build-out}
\begin{itemize}
\item \textbf{Trial Balance:} Implement Controls \& Commentary Engine (3$\sigma$ anomaly logic). See \cref{tb:ch:ai,tb:ch:controls-engine}.
\item \textbf{Arabic:} Implement multi-country VAT Engine (Saudi, UAE, Bahrain, Oman). See \cref{ar:ch:vat,ar:ch:vat-impl}.
\item \textbf{Simula:} Implement Taxonomy Generation \& Agentic Data Synthesis. See \cref{sim:chap:taxonomy,sim:chap:generator}.
\item \textbf{Trial Balance:} Realise 28-state BPMN workflow (\texttt{tb-review-glo}). See \cref{tb:ch:workflow}.
\end{itemize}

\subsection*{Wave 3 --- Hardening}
\begin{itemize}
\item \textbf{Regulations:} Implement Per-Request Identity Attribution. See \cref{reg:ch:identity}.
\item \textbf{Trial Balance:} Deliver Human Review Dashboard. See \cref{tb:ch:review-dashboard}.
\item \textbf{Simula:} Deliver test fixtures and testing at 85\% coverage. See \cref{sim:ch:test-fixtures,sim:ch:testing}.
\item \textbf{Regulations:} Wire Emergent-Capability monitoring. See \cref{reg:ch:monitoring}.
\end{itemize}

\section{Governance Obligations (MAS MGF)}
The regulations domain wraps all service calls at runtime. Agents must implement the following hooks in every domain:
\begin{itemize}
\item \textbf{Pillar 1:} Tool allow-list enforcement (deny-by-default).
\item \textbf{Pillar 2:} Traceability via the identity attribution envelope.
\item \textbf{Pillar 3:} Metric emission to the monitoring sink.
\end{itemize}
Detailed pillar requirements are defined in \cref{reg:ch:mgf}.

\section{FMEA: Failure Mode and Effects Analysis for AI}
\label{sec:fmea-ai}

The following analysis identifies critical AI failure modes within the Arabic AP pipeline and their deterministic mitigations.

\begin{table}[htbp]
\centering
\footnotesize
\caption{FMEA for Arabic AP AI Pipeline.}
\label{tab:arabic-fmea}
\begin{tabularx}{\textwidth}{@{}l X X X@{}}
\toprule
\textbf{Failure Mode} & \textbf{Symptom} & \textbf{Impact} & \textbf{Technical Mitigation} \\
\midrule
OCR Character Drift & Misreading '5' as '6' in currency fields. & High (Financial) & Cross-field arithmetic validation (Net + VAT == Gross). \\
Translation Ambiguity & 'Invoice' translated as 'Note' (misclassification). & Medium (Routing) & Metadata-tagging based on ZATCA QR-code extraction. \\
LLM Field Extraction & Extracting 'Address' into 'Vendor Name'. & Medium (Data Qual.) & Confidence-score threshold check per field. \\
\bottomrule
\end{tabularx}
\end{table}

\section{This domain in context}
\label{sec:agent-domain-context}
Arabic AP is the highest-volume transactional workstream among the
four domains. Its primary persona is the AP Clerk, whose working view
is the approval queue surfaced through the shared workspace shell and
gated by AP entitlements. The domain's entry points are document
intake (email, SFTP, upload) producing raw invoice artefacts; its
exit points are (a)~a posted AP document in SAP S/4HANA with a signed
approval record, and (b)~a VAT compliance checklist emitted to the
audit store. Every artefact is validated against the Arabic-owned
schemas in \texttt{structured/schema/} before it leaves the service.
The domain depends on the horizontal regulations wrapper for runtime
governance and on Simula for the training corpus used by the OCR,
field-extraction, and translation models described in
chapter~\ref{ch:ai}.

\section{Schema contribution}
\label{sec:agent-schema-contrib}
Arabic AP owns the schemas catalogued in chapter~\ref{ch:schema}. The
canonical set is the invoice record (see~\ref{tab:invoice-fields}),
the VAT compliance checklist (see~\ref{tab:checklist-fields}), and the
payment instruction record. These schemas live under
\texttt{docs/latex/specs/arabic/} as normative text and under
\texttt{structured/schema/} as the JSON Schema artefacts consumed at
runtime. Registration in the unified registry is additive: each
schema is listed in
\texttt{src/intelligence/ai-core-pal/data\_products/registry.yaml}
with its URI, owning domain (\texttt{arabic}), security class
(\texttt{confidential}), and LLM routing
(\texttt{vllm-only}), following the ODPS~4.1 pattern. Simula consumes
these schemas directly as inputs to its extraction pipeline; it does
not transform them. The \texttt{requirement.schema.json} files owned
by TB and regulations remain distinct from anything Arabic owns and
are not merged with it.

\section{Persona and MCP allow-list}
\label{sec:agent-persona}
\begin{itemize}
\item \textbf{Persona id:} \texttt{ap-clerk}.
\item \textbf{Roles:} invoice intake, VAT review, approval
progression, exception handling.
\item \textbf{Entitlements:} read/write on the Arabic invoice
record; read on the VAT rule set; advance transitions on the
\texttt{po-eproc} and \texttt{direct-ap} state machines defined
in chapter~\ref{ch:ap}; evidence-write on the control catalogue
in~\ref{tab:controls}.
\item \textbf{UI route:} \texttt{/ap/\{tenant\}/queue} on the shared
workspace shell; surfaced only to callers whose entitlement set
includes \texttt{ap-clerk}.
\item \textbf{MCP tool allow-list:}
\texttt{invoice.intake}, \texttt{invoice.ocr},
\texttt{invoice.extract-fields},
\texttt{invoice.translate-ar-en}, \texttt{vat.validate},
\texttt{vat.calculate}, \texttt{ap.workflow.advance},
\texttt{ap.workflow.park}. All other MCP tools are denied by
default per MAS MGF pillar~1.
\end{itemize}

\section{Implementation tasks for this domain}
\label{sec:agent-tasks}
The tasks below are intent-based: each task states the outcome a
future agent must achieve and the location in the existing codebase
where the work belongs. None of them contains code; each task binds
to the Arabic chapters that define the required behaviour.

\subsection*{Task A1 --- OCR and translation calibration}
\textbf{Intent.} Bring Arabic OCR and Arabic--English translation up
to the precision/recall targets in chapter~\ref{ch:ocr-pilot}.

\begin{adrbox}{Why Tesseract + Custom Post-Processor?}
We chose Tesseract 5.0 for the baseline due to local-hosting requirements, but added a custom 'Arabic-Context' post-processor to handle the right-to-left layout drift in scanned PDF tables.
\end{adrbox}

\textbf{In-scope.} Running the pilot sample defined
in~\ref{tab:pilot-sample}; publishing per-field thresholds for
accept/review/reject; wiring the thresholds into the pipeline stages
of chapter~\ref{ch:ai}.
\textbf{Out-of-scope.} Choosing a new OCR engine; altering the
canonical invoice schema.
\textbf{Paths to extend.} OCR and extraction code colocated with the
pipeline stages referenced in chapter~\ref{ch:ai}; fixture store
referenced in chapter~\ref{ch:testing}.
\textbf{Definition of done.} Pilot report committed; thresholds
checked in as configuration; the stage-3 field extractor rejects
records below threshold and routes them to human review.

\textbf{Autonomous Verification (Agent Logic):}
\begin{lstlisting}[language=jsonx]
// Verify Task A1
{
"test_fixture": "tests/fixtures/arabic/invoice_sample_01.png",
"expected_outcome": "ocr_confidence > 0.88",
"validation_command": "make test-arabic-ocr"
}
\end{lstlisting}

\textbf{Verification.} \texttt{make spec-arabic} builds cleanly and
the Arabic test suite described in chapter~\ref{ch:testing} passes
end-to-end on the calibration fixtures.

\subsection*{Task A2 --- Multi-country VAT engine}
\textbf{Intent.} Implement the VAT compliance engine specified in
chapter~\ref{ch:vat} and detailed in chapter~\ref{ch:vat-impl}.
\textbf{In-scope.} Per-country rule modules for Saudi Arabia (ZATCA),
UAE, Bahrain (NBR), and Oman; registration validation, calculation,
activity lookup, and checklist emission; the canonical decision
function in~\ref{fig:ksa-decision}.
\textbf{Out-of-scope.} VAT rules for countries outside the in-scope
list in~\ref{tab:countries}.
\textbf{Paths to extend.} The VAT module layout referenced
in~\ref{fig:vat-arch}; API contracts in chapter~\ref{ch:vat-impl}.
\textbf{Definition of done.} Every rule cited
in~\ref{tab:ksa-rules} has a callable check; the checklist record is
emitted per invoice and validates against
\texttt{vat-checklist.schema.json}.
\textbf{Verification.} Unit tests for each country module (see
chapter~\ref{ch:testing}) pass; the VAT calculation API contract
sample in chapter~\ref{ch:vat-impl} returns the expected response.

\subsection*{Task A3 --- AP approval workflow}
\textbf{Intent.} Realise the AP workflow and approval engine in
chapter~\ref{ch:ap}, including the \texttt{po-eproc} and
\texttt{direct-ap} state machines.
\textbf{In-scope.} Workflow-selection logic in~\ref{fig:workflow-select};
state transitions in~\ref{fig:po-eproc}; park, reroute, and escalate
actions with evidence records.
\textbf{Out-of-scope.} SAP S/4HANA AP posting internals;
non-Bahrain delegations.
\textbf{Paths to extend.} Workflow engine hosted in the gateway
service; evidence store colocated with the control catalogue.
\textbf{Definition of done.} Each state in the state-machine diagram
is reachable from a deterministic predecessor; transitions are audit
logged with the persona id of the acting user; the golden regression
set in chapter~\ref{ch:controls} executes without human intervention.
\textbf{Verification.} Workflow unit tests pass; the approval queue
UI surfaces the correct entitlements for \texttt{ap-clerk}.

\subsection*{Task A4 --- Bahrain DOI seven controls}
\textbf{Intent.} Evidence the seven Bahrain Delegation of Invoicing
controls required in chapter~\ref{ch:controls}.
\textbf{In-scope.} Control enforcement map in~\ref{fig:controls-map};
the seven rows of the control catalogue in~\ref{tab:controls} that
concern the Bahrain DOI; evidence emission per control.
\textbf{Out-of-scope.} Introducing controls for other delegations.
\textbf{Paths to extend.} Control catalogue file referenced in
chapter~\ref{ch:controls}; evidence store.
\textbf{Definition of done.} Running the workflow on the Bahrain
golden sample triggers every control; each control writes an
evidence record conforming to the Arabic schema set.
\textbf{Verification.} Traceability matrix in chapter~\ref{ch:controls}
is complete; the testing chapter's acceptance suite passes.

\subsection*{Task A5 --- Testing and golden regression set}
\textbf{Intent.} Deliver the testing specification in
chapter~\ref{ch:testing} at the coverage targets
in~\ref{tab:arabic-coverage}.
\textbf{In-scope.} Invoice and OCR fixtures; unit, integration, and
regression suites; CI wiring.
\textbf{Out-of-scope.} Performance benchmarking infrastructure beyond
what chapter~\ref{ch:testing} defines.
\textbf{Paths to extend.} Test directory colocated with the Arabic
service code; fixture store used by Simula as training input.
\textbf{Definition of done.} Coverage meets the thresholds
in~\ref{tab:arabic-coverage}; CI gates block regressions; the
regression golden set is reproducible from fixture seeds.
\textbf{Verification.} \texttt{make spec-arabic} remains green; the
Arabic test entry points listed in chapter~\ref{ch:testing} all run.


\section{Governance obligations applied to this domain}
\label{sec:agent-governance}
The regulations domain wraps Arabic AP at runtime. Four obligations
apply without exception; each is stated as a concrete hook in the
Arabic codebase.

\begin{itemize}
\item \textbf{MAS MGF four pillars} (see regulations
chapter~3 in \texttt{docs/latex/specs/regulations/chapters/03-mgf-framework.tex}).
Pillar~1 (assess and bound): every Arabic MCP tool is declared in
the allow-list in \S\ref{sec:agent-persona}, deny-by-default.
Pillar~2 (design and deploy): the pipeline stages in
chapter~\ref{ch:ai} are traceable end-to-end via the
per-request identity attribution described below. Pillar~3
(operate and monitor) and pillar~4 (govern and disclose) are
satisfied by emitting the evidence records produced by the
Arabic control catalogue (chapter~\ref{ch:controls}) into the
shared governance store.
\item \textbf{Identity attribution} (regulations chapter~10,
\texttt{10-identity-attribution.tex}). Every MCP call and every
workflow transition carries the persona id, agent id, request
id, and tenant id. The gateway is the attachment point; Arabic
code propagates the identity envelope without mutating it.
\item \textbf{Emergent-capability monitoring} (regulations
chapter~8, \texttt{08-monitoring-spec.tex}). OCR, extraction,
translation, and VAT calculation are instrumented with the
monitoring metrics specified in that chapter; metrics flow to
the monitoring sink defined there.
\item \textbf{Conformance tooling} (regulations chapter~6,
\texttt{06-conformance-tooling.tex}). The Arabic service exposes
AI Verify 2.0 plug-in hooks and Moonshot CI/CD 1.1.0 benchmark
entry points as required by that chapter.
\end{itemize}

\section{Cross-spec integration --- full detail}
\label{sec:agent-integration}

\subsection{With Arabic AP}
Arabic AP exports the invoice record, VAT compliance checklist, and
payment instruction schemas under \texttt{structured/schema/}
(chapter~\ref{ch:schema}); registers the MCP tools enumerated in
\S\ref{sec:agent-persona}; and applies the four governance hooks in
\S\ref{sec:agent-governance} to every call it serves. External
consumers (TB, regulations, Simula) see Arabic only through these
contracts; there are no private back-channels.

\subsection{With Trial Balance}
Arabic AP produces posted AP records and VAT checklists that feed
the TB close. TB consumes them as evidence inputs to its
commentary engine, not as schema merges: the Arabic invoice schema
and the TB control record remain distinct URIs in the unified
registry. Integration point: TB's ingestion stage calls the
Arabic MCP tool \texttt{invoice.extract-fields} with read-only
entitlements; the response is validated against
\texttt{invoice.schema.json} before TB accepts it.

\subsection{With Regulations}
Regulations wraps Arabic at runtime. Every Arabic MCP call is
brokered through the regulations governance gate, which applies
pillars~1--4 of MAS MGF, attaches the identity envelope, and routes
emergent-capability metrics to the monitoring store. The Arabic
service code does not call regulations directly; the gateway layer
enforces this. Conformance tests defined in regulations
chapter~5 (empirical evidence) and chapter~6 (conformance tooling)
execute against the Arabic service as part of the platform CI gate.

\subsection{With Simula}
Simula consumes Arabic's schemas as training input. The contract is
schema-in, \texttt{TrainingExample}-out: Simula reads
\texttt{invoice.schema.json}, \texttt{vat-checklist.schema.json}, and
the Arabic fixture store; it produces \texttt{TrainingExample}
records. Simula does not call Arabic MCP tools at runtime and does
not share state with the Arabic workflow. No cross-domain
operational envelope is introduced.

\section{Wave assignment and sequencing}
\label{sec:agent-waves}
Implementation runs in three waves.

\begin{itemize}
\item \textbf{Wave 1 --- foundations.} Unified schema registry
entry; persona-gated surface for the AP Clerk; identity
attribution envelope flowing through the gateway. Prerequisites:
the regulations chapter~10 identity schema must be frozen.
\item \textbf{Wave 2 --- Arabic core build-out.} Tasks A1--A4 above.
Prerequisites: Wave~1 complete; Simula training fixtures
available from Simula Wave~1 (extraction pipeline).
\item \textbf{Wave 3 --- hardening and conformance.} Task A5
and full conformance-tooling hook-up. Prerequisites: Wave~2
complete; regulations Wave~2 conformance tooling integrated;
monitoring sink live.
\end{itemize}

Arabic occupies Wave~2 for its core build-out and Wave~3 for its
testing and conformance surfaces.

\section{Acceptance criteria (platform-wide)}
\label{sec:agent-acceptance}
When all four chapter~12 instruction sets are implemented, the
platform satisfies the following invariants:

\begin{enumerate}
\item \textbf{One instance.} A single codebase runs a single
logical instance; no domain ships as a separate product.
\item \textbf{Persona-gated surfaces.} The four personas each
reach a tailored UI on the shared shell; no route, MCP tool, or
data product is visible outside its persona's entitlement set.
\item \textbf{Unified registry.} Every schema in use is listed in
the ODPS~4.1 registry under its owning domain; validators
resolve \texttt{\$ref}s through the unified registry only.
\item \textbf{No schema merges.} The two
\texttt{requirement.schema.json} files in TB and regulations
remain distinct; no cross-domain composite schemas are produced.
\item \textbf{No cross-domain operational envelope.} Runtime
orchestration stays within each domain's service; only Simula
consumes other domains' outputs, and only as training input.
\item \textbf{Governance wraps all three non-regulations domains.}
Arabic, TB, and Simula service calls pass through the
regulations runtime wrapper without exception.
\end{enumerate}

\section{References}
\label{sec:agent-references}
Within-spec references:
chapter~\ref{ch:overview} (overview),
chapter~\ref{ch:schema} (data schema),
chapter~\ref{ch:ai} (AI pipeline),
chapter~\ref{ch:vat} (VAT compliance engine),
chapter~\ref{ch:ap} (AP workflow),
chapter~\ref{ch:controls} (controls and testing),
chapter~\ref{ch:references} (references),
chapter~\ref{ch:vat-impl} (VAT implementation),
chapter~\ref{ch:testing} (testing specification),
chapter~\ref{ch:ocr-pilot} (OCR pilot),
chapter~\ref{ch:ap-roadmap} (project roadmap).

Peer chapter~12 instruction sets for the other three domains (each
is self-contained; consult directly when operating in that domain):
\begin{itemize}
\item Trial Balance:
\texttt{docs/latex/specs/tb/chapters/12-implementation-instructions.tex}.
\item Regulations:
\texttt{docs/latex/specs/regulations/chapters/12-implementation-instructions.tex}.
\item Simula:
\texttt{docs/latex/specs/simula/chapters/12-implementation-instructions.tex}.
\end{itemize}

Shared platform artefacts:
\begin{itemize}
\item ODPS~4.1 registry:
\texttt{src/intelligence/ai-core-pal/data\_products/registry.yaml}.
\item Schema pipeline and resolver:
\texttt{src/intelligence/schema\_pipeline/}.
\item Training pipeline (Simula):
\texttt{src/training/pipeline/}.
\item Shared workspace shell:
\texttt{generativeUI/training-webcomponents-ngx}.
\end{itemize}

% --- Incorporated Universal Chapters ---
\chapter{Global Wave-Dependency Manifest}
\label{ch:global-dependencies}

\section{Critical Path Visualization (Waves 1--3)}
The following chart defines the normative sequencing for implementation agents. A task in a later wave cannot be marked \textsc{complete} until its prerequisites in an earlier wave are validated.

\begin{figure}[htbp]
\centering
\begin{tikzpicture}[
node distance=1.5cm,
auto,
wavebox/.style={rectangle, draw, rounded corners, fill=sapmidnight!10, text width=10em, text centered, minimum height=3em, font=\scriptsize\bfseries},
dep/.style={->, >=latex, thick, sapblue},
critical/.style={->, >=latex, very thick, sapamber}
]
% --- Wave 1: Foundations ---
\node (r1) [wavebox, fill=sapgrey!20] {\textcolor{sapmidnight}{REG-W1: Identity \& Audit Schema}};
\node (s1) [wavebox, below of=r1, yshift=-0.5cm] {SIM-W1: HANA Extraction Pipeline};
\node (a1) [wavebox, below of=s1, yshift=-0.5cm] {ARA-W1: Invoice Intake \& OCR};
\node (t1) [wavebox, below of=a1, yshift=-0.5cm] {TB-W1: Identity Wiring};

% --- Wave 2: Core Build-out ---
\node (r2) [wavebox, right of=r1, xshift=4.5cm, fill=sapmidnight!20] {\textcolor{sapmidnight}{REG-W2: Governance Runtime}};
\node (s2) [wavebox, right of=s1, xshift=4.5cm] {SIM-W2: Taxonomy \& Synthesis};
\node (a2) [wavebox, right of=a1, xshift=4.5cm] {ARA-W2: Multi-Country VAT};
\node (t2) [wavebox, right of=t1, xshift=4.5cm] {TB-W2: Controls \& Workflow};

% --- Wave 3: Hardening ---
\node (r3) [wavebox, right of=r2, xshift=4.5cm, fill=sapmidnight!30] {\textcolor{sapmidnight}{REG-W3: Conformance Suite}};
\node (s3) [wavebox, right of=s2, xshift=4.5cm] {SIM-W3: 85\% Coverage Test};
\node (a3) [wavebox, right of=a2, xshift=4.5cm] {ARA-W3: Calibration Pilot};
\node (t3) [wavebox, right of=t2, xshift=4.5cm] {TB-W3: HITL Dashboard};

% --- Cross-Domain Dependencies (The Critical Path) ---
\draw [critical] (r1) -- node [above, font=\tiny, align=center] {Identity Envelope\\Invariants} (a1.north east);
\draw [critical] (r1) -- (t1.north east);
\draw [dep] (s1.east) -- node [above, font=\tiny, xshift=-0.5cm] {Training Fixtures} (t2.west);
\draw [dep] (s1.east) -- (a2.west);
\draw [dep] (r2.south) -- node [left, font=\tiny, yshift=0.3cm] {Policy Gate} (a2.north);
\draw [dep] (t2.east) -- node [above, font=\tiny] {Audit Evidence} (r3.west);

% --- Wave Transitions ---
\draw [->, dashed, sapgrey] (r1.east) -- (r2.west);
\draw [->, dashed, sapgrey] (r2.east) -- (r3.west);

% --- Legend ---
\node[draw, sapamber, very thick, label=right:\scriptsize Critical Prerequisite] at (0,-7.5) {};
\node[draw, sapblue, thick, label=right:\scriptsize Data Dependency, xshift=4cm] at (0,-7.5) {};

\end{tikzpicture}
\caption{Cross-specification dependency graph: The blueprint for platform rollout.}
\end{figure}

\section{Inter-Domain Prerequisite Matrix}

\begin{table}[htbp]
\centering
\footnotesize
\caption{Normative Cross-Domain Prerequisites.}
\label{tab:prerequisites}
\begin{tabularx}{\textwidth}{@{}l l l X@{}}
\toprule
\textbf{Target Task} & \textbf{Prerequisite} & \textbf{Ref ID} & \textbf{Rationale} \\
\midrule
\rowcolor{sapmidnight!5}
\textbf{ARA-W1} & REG-W1 (Identity) & \texttt{REG-IDENTITY-01} & Invoice intake must be audit-traced from Day 1. \\
\textbf{TB-W2}  & SIM-W1 (Extraction) & \texttt{SIM-EXTRACT-01} & Anomaly detection requires HANA schema metadata. \\
\rowcolor{sapmidnight!5}
\textbf{REG-W2} & SIM-W1 (Registry) & \texttt{SIM-REGISTRY-01} & Governance monitoring relies on the schema registry. \\
\textbf{ARA-W2} & REG-W2 (Pillar 1) & \texttt{REG-MGF-P1} & VAT release is an HITL gate governed by MAS Pillar 1. \\
\rowcolor{sapmidnight!5}
\textbf{TB-W3}  & ARA-W2 (Posted Invoices) & \texttt{ARA-AP-POST} & Review dashboard depends on posted AP data from Arabic domain. \\
\bottomrule
\end{tabularx}
\end{table}

\begin{adrbox}{Dependency Enforcement}
Implementation agents are strictly forbidden from "Fast-Forwarding" into Wave 2 if Wave 1 critical prerequisites (marked in \textcolor{sapamber}{\textbf{Gold}}) are in a \textsc{gap} or \textsc{partial} state.
\end{adrbox}



\end{document}
